%--- iliad ---
% cluster: B
% title: "Singular Learning Theory"
%--- end ---
\documentclass[11pt]{article}
% % %
% Packages

\usepackage[utf8]{inputenc}
\usepackage[T1]{fontenc}
\usepackage{mathtools}
\usepackage{fancyhdr}
\usepackage[round,compress]{natbib}

% iliad.sty: local copy first (Overleaf), else the shared ../iliad.sty.
% Provides exercise/solution/definition/theorem/lemma/proposition/corollary/
% example/remark/callout, the \ifsolutions toggle, and loads
% amsmath+amssymb+amsthm+enumitem+graphicx+tcolorbox+hyperref+cleveref.
\IfFileExists{iliad.sty}{\usepackage[boxes]{iliad}}{\usepackage[boxes]{../iliad}}

% Restore the original link colour (iliad defaults to plain blue).
\hypersetup{colorlinks, allcolors=blue!60!black}

% % %
% Configuration

% Table of contents
\setcounter{tocdepth}{2}

% Indentation
\setlength{\parindent}{0pt}
\setlength{\parskip}{1ex}

% Bibstyle (natbib-compatible; far.bst ships in this folder)
\bibliographystyle{far}

% Custom copyright footer
\fancypagestyle{copyright}{
    \fancyhf{}
    \renewcommand{\headrulewidth}{0pt}
    \cfoot{\copyright{} Kai Ogden, Matthew Farrugia-Roberts, Zach Furman, 2026. All rights reserved.}
}

% % % 
% Notation
\newcommand{\R}{\mathbb{R}}
\newcommand{\N}{\mathbb{N}}
\newcommand{\E}{\mathbb{E}}
\newcommand{\KL}{\mathrm{KL}}
\newcommand{\tr}{\mathrm{tr}}
\newcommand{\rank}{\mathrm{rank}}
\newcommand{\im}{\mathrm{im}}
\newcommand{\sign}{\mathrm{sign}}
\newcommand{\relu}{\mathrm{relu}}
\DeclareMathOperator*{\argmin}{arg\,min}
\DeclareMathOperator*{\argmax}{arg\,max}
\newcommand{\X}{\mathcal{X}}
\newcommand{\Y}{\mathcal{Y}}
\newcommand{\W}{\mathcal{W}}
\newcommand{\F}{\mathcal{F}}
\newcommand{\D}{\mathcal{D}}
\newcommand{\C}{\mathcal{C}}
\newcommand{\U}{\mathcal{U}}
\newcommand{\V}{\mathcal{V}}



% % % 
% Frontmatter


\title{
    \textbf{Degeneracy in Deep Learning}\\
    An Invitation to Singular Learning Theory
}
\author{
    \authorname{Kai Ogden}\\
    \affiliation{University of Oxford}
    \and
    \authorname{Matthew Farrugia-Roberts}\\
    \affiliation{University of Oxford}
    \and
    \authorname{Zach Furman}\\
    \affiliation{The University of Melbourne}
}
\date{%
    Pilot, Spring 2026
}

\begin{document}

% % % 
% Title/TOC page

\maketitle

\thispagestyle{copyright}

\bgroup
% disable title
\renewcommand\section[2]{}
% disable spacing
\setlength{\parskip}{0ex}
% print the toc
\tableofcontents
\egroup


% % % 
% Epigraph page

\clearpage
\thispagestyle{empty}
\null
\vfill
\begin{quote}
    \centering
    \textit{Thou shalt have the power to degenerate}\\
    \qquad\qquad\qquad---\textit{On the Dignity of Man,} 1496
\end{quote}
\vfill
\vfill


% % % 
% Begin body
\clearpage

\begin{summary}
Singular learning theory (SLT) places \emph{degeneracy} as a core part of
understanding how neural networks learn.
We cover the parameter-function map, the meaning of degeneracy
through the local learning coefficient, to Watanabe's free energy formula and
Bayesian phase transitions.
\end{summary}

\section*{Introduction}
\label{sec:intro}
\addcontentsline{toc}{section}{\nameref{sec:intro}}

\emph{Singular learning theory (SLT)} is a theory of learning that accounts for
\emph{degeneracy} in neural networks. What is degeneracy? Let us first appeal
to the Cambridge dictionary, documenting some features of typical uses of the
word within mathematics:
\begin{quote}
    \emph{Degenerate} (of an equation, curve, line, etc.) unusual or
    complicated in some way compared to other equations, curves, lines, etc.\
    of a similar type, especially because a variable or parameter is zero.
\end{quote}
On the other hand, we should recall that deep learning has as many roots in
neuroscience as in mathematics. And as they say, ``neural networks are
\emph{grown.}'' Perhaps we should also hear the biologist's definition:
\begin{quote}
    \emph{Degenerate} (of an organism, chromosome, etc.) simpler than a form
    that previously existed because of no longer having a particular structure.
\end{quote}

Actually, \emph{both} of these definitions are apt. Within a typical neural
network architecture, certain weight vectors correspond to neural networks from
simpler architectures (often due to certain weights being zero). Structurally,
these neural networks are simpler in form than their neighbours.
Mathematically, they complicate the relationship between the neural network's
parameter space and the resulting space of functions. In turn, learning in
neural networks is substantially richer than learning in classical statistical
models.

Singular learning theory is a framework for understanding learning that places
degeneracy at the centre. The framework leverages powerful tools from algebraic
geometry to understand, characterise, and resolve degeneracies, revealing their
impact on learning. Unfortunately, the price of this mathematical power is that
following the research literature without a deep background in pure mathematics
is challenging.

That is where this tutorial comes in. We aim to distil the essence of singular
learning theory---the notion of degeneracy and its role in learning---laying
bare through simple examples and structured exercises the core intuitions
driving research in the field. We believe understanding this essence requires
no more than undergraduate-level mathematics. Moreover, to anyone armed with
these core intuitions, understanding and contributing to research on the
relationship between degeneracy and deep learning is within reach.

\paragraph{Contents.} Specifically, the tutorial comprises the following four
technical sections, each with definitions and guided exercises.
\begin{enumerate}
    \item
        ``\nameref{sec:prelims}'' reviews core deep learning concepts, defining
        our notation and some running examples. The central concept is that of
        the parameter--function map.
    \item
        ``\nameref{sec:degeneracy}'' defines degeneracy of parameter--function
        maps and relates it to symmetries, singularities of the Fisher
        information, and degenerate critical points of the loss landscape.
    \item
        ``\nameref{sec:llc}'' derives the local learning coefficient (LLC) as a
        quantitative measure of the degree of degeneracy of a critical point
        via a volume scaling approach, and surveys several other perspectives
        on the quantity.
    \item
        ``\nameref{sec:fef}'' presents Watanabe's free energy formula
        connecting the LLC to Bayesian learning and discusses its implications.
\end{enumerate}
We conclude with \cref{sec:readings}, providing pointers for further study and
surveying the literature on SLT and deep learning. 
\begin{solutionsonly}
\Cref{apx:solutions}
provides worked solutions for all exercises.
\end{solutionsonly}
\paragraph{Prerequisites.}

The following is an indicative list of mathematical concepts that will be
helpful for reading the tutorial and completing the exercises.

\begin{itemize}
    \item
        \emph{Linear algebra:} vectors, matrices, rank, orthogonal matrices,
        rank--nullity, positive definiteness, eigenvalues, spectral decomposition.
    \item
        \emph{Calculus:} partial derivatives, gradient, directional derivative,
        chain rule, Hessian, second-order Taylor expansion and remainder.
    \item
        \emph{Integration and analysis:} multivariate integrals, change of
        variables, volume in $\R^d$, asymptotic notation (big-$O$, little-$o$),
        computing basic limits and integrals.
    \item
        \emph{Probability:} probability simplex $\Delta(\cdot)$, conditional
        probability, probability density functions, independence, expectation,
        Bayes' rule, Gaussians, law of large numbers.
\end{itemize}

\Cref{sec:prelims} reviews the basic framework of deep learning as parametric
function approximation or statistical inference (parameter--function maps, deep
linear networks, multi-layer perceptrons, loss functions, likelihood) along
with Bayesian inference (prior, posterior, partition function, Bayesian free
energy).

Readers with more advanced backgrounds may appreciate occasional references to
topics from algebraic geometry, fractal geometry, or statistical physics.
However, these references are tangential and readers without these backgrounds
can safely skip them.

\paragraph{Fast-track.}

To paraphrase Euclid, there is no royal road to algebro-geometric learning
theory. However, it is possible to get a bird's-eye view and the most important
intuitions in a comparably short time. If this is what you are looking for, we
recommend the following route through the tutorial. Assuming you are already
somewhat comfortable with deep learning and Bayesian inference, skip
\cref{sec:prelims}, and refer back only as needed. Then, proceed as follows.
\begin{enumerate}
    \item
        To understand parameter--function map versus loss landscape degeneracy:
        Read \cref{sec:degeneracy:definition} and complete
        \cref{ex:univariate-degeneracy,ex:cubic-degeneracy}.
        Complete either \cref{ex:dln-degeneracy} or \ref{ex:mlp-degeneracy}.
        Read \cref{sec:degeneracy:hessian} and complete your choice of
            Exercises~\ref{ex:loss-landscape-degeneracy} and/or \ref{ex:fim-hessian}.
    \item
        To understand the local learning coefficient via volume scaling:
        Read \cref{sec:llc:volume-scaling} and complete
        \cref{ex:llc-regular-case,ex:quad-valley-scaling,ex:calculating-llc,ex:cubic-parameterisation}.
        Read \cref{sec:llc:dln} and \cref{ex:dln-lc}.
    \item
        To understand the relation between degeneracy and learning in the
        Bayesian case:
        Read all of \cref{sec:fef} (it is shorter).
        Complete \cref{ex:phase-transitions}.
\end{enumerate}
Once you are done, we hope you will consider taking the scenic route some other
time.

\paragraph{Acknowledgements.}

This tutorial was developed for the April 2026 Iliad Intensive.
Some exercises were adapted from \citet{Furman2024}.

\clearpage
\section{Preliminaries}
\label{sec:prelims}

In this section, we introduce basic terminology and notation for supervised
deep learning and supervised Bayesian deep learning, along with some example
neural network architectures and statistical models that we will repeatedly
study throughout the tutorial.

\subsection{Neural networks and parameter--function maps}

Let $\X$ denote a space of network inputs (e.g., a space of images or token
sequences encoded as vectors), and let $\Y$ denote a space of outputs (e.g., a
space of numerical scores, class distributions, or next-token distributions).

Informally, a neural network architecture is a specification for how various
\textbf{parameters} (e.g., neuron connection strengths / weights or neuron
activation thresholds / biases) combine with input signals and each-other to
describe a function from $\X$ to $\Y$.

Formally, a neural network architecture essentially comprises a
\textbf{parameter--function map}
    $\Phi : \W \to \F$, where
$\W \subseteq \R^d$ is a $d$-dimensional \textbf{parameter space} and 
$\F \subseteq \Y^\X$ is some \textbf{hypothesis class} of functions from $\X$
to $\Y$.
%
It is sometimes convenient to refer to the function $\Phi(w) : \X \to \Y$ as
$f_w : \X \to \Y$.

Throughout the rest of this tutorial, we will study many different neural
network architectures (parameter--function maps). The rest of this section
explores some generally useful examples and some terminological points. To
begin with, we have the following remark.

\begin{remark}[Parameters]
    Note that the term ``parameter'' has two senses:
    \begin{enumerate}
        \item
            A ``parameter'' is an individual dimension of parameter space
            (e.g., ``initialise this parameter to the value $0.0$,'' or ``this
            neural network has billions of parameters''). A loose synonym in
            this case is ``weight.''
        \item
            A ``parameter'' is also a particular point in the parameter space
            (e.g., ``this parameter is a local minimum of the loss function,''
            or ``the zero locus of this loss function contains a continuum of
            parameters''). A loose synonym in this case is ``weight vector.''
    \end{enumerate}
    Both senses are in common usage in the literature as in this tutorial.
\end{remark}

The simplest neural network architecture models a single ``neuron'' with $d$
inputs $x_1, \ldots x_d$. Each input $x_i$ is multiplied by an incoming weight
$w_i$ to produce the neuron's output. We formalise this case in
\cref{eg:neuron}.

\begin{example}[A linear neuron]\label{eg:neuron}
    Let $\X = \R^d$ for some positive integer $d$ and let $\Y = \R$. Define a
    parameter space $\W = \R^d$. Define a parameter--function map that maps a
    column vector $w \in \W$ to the function $f_w : \R^d \to \R$ such that for
    $x \in \R^d$,
    \begin{equation}
        f_w(x)
        = w^\top x
        = \sum_{i=1}^d w_i x_i.
    \end{equation}
\end{example}

Let us generalise this basic neural network in three ways: to add multiple
outputs, ``depth,'' and non-linearity. First, we can generalise from a scalar
output to a vector output as follows.

\begin{example}[Multi-linear neural network]\label{eg:linear}
    Let $\X = \Y = \R^m$ for some positive integer $m$.
    Define a parameter space $\W = \R^{m^2}$. Let each parameter $w \in \W$
    encode an $m \times m$ matrix $W \in \R^{m \times m}$. The
    parameter--function map transforms each parameter $w \in \W$ to a function
    $f_w : \R^m \to \R^m$ such that for $x \in \R^m$,
    \begin{equation}
        f_w(x) = W x.
    \end{equation}
\end{example}

If we take each row of $W$ to encode the weight vector of a linear neuron
(\cref{eg:neuron}), we see that we have produced $m$ outputs by stacking $m$
independent linear neurons together. Such a group of neurons is called a
\textbf{layer}.

\begin{remark}[Encodings]
    In \cref{eg:linear}, the parameter space is formally $\W = \R^{m^2}$, but
    we find it convenient to identify each parameter vector $w \in \R^{m^2}$
    with the $m \times m$ matrix $W$ it encodes. We would thus write $f_W$ in
    place of $f_w$. More generally, whenever the parameters of an architecture
    naturally decompose into named matrices or vectors, we index the
    parameter--function map by these structured objects directly rather than by
    the underlying vector $w$. The remaining examples in this section
    demonstrate this convention, and we continue to use it throughout this
    tutorial.
\end{remark}

Next, let's add \emph{depth} to our neural network by composing two layers
together.

\begin{example}[Deep linear network; DLN]\label{eg:dln}
    Let $\X = \Y = \R^m$ for some positive integer $m$.
    Let $h$ be a positive integer and define the parameter space $\W =
    \R^{2mh}$, with each parameter encoding a pair of matrices
        $A \in \R^{h \times m}$
    and $B \in \R^{m \times h}$.
    The parameter--function map sends $(A, B)$ to $f_{A,B} : \R^m \to \R^m$
    such that for $x \in \R^m$,
    \begin{equation}
        f_{A,B}(x) = B A x.
    \end{equation}
\end{example}

The above architecture is called a two-layer \textbf{deep linear network
(DLN)}. The architecture can be interpreted as composing two multi-linear
neural networks. The vector of outputs of the neurons of the first network
becomes the vector of inputs to the second network. The intermediate output
vectors are called \textbf{activations}.

Note that if $h \geq m$, the two-layer DLN architecture indexes the same
hypothesis class as in \cref{eg:linear}, that of linear transforms on $\R^m$.
If $h < m$, the hypothesis class includes only transforms with rank up to $h$.

Finally, we'll add \emph{non-linearity} between pairs of layers. To do so, we
introduce a non-linear scalar function $\sigma : \R \to \R$, called an
\emph{activation function}, to transform the output of each intermediate
neuron. Common examples of activation functions, which we will study in this
tutorial, include the following:
\begin{enumerate}
    \item
        The \textbf{hyperbolic tangent} function
            $\displaystyle \tanh(z) = \frac{e^z-e^{-z}}{e^z + e^{-z}}$.
    \item
        The \textbf{rectified linear unit (ReLU)} function $\relu(z) = \max\{0, z\}$.
\end{enumerate}
This results in the following non-linear neural network architecture.

\begin{example}[Multi-layer perceptron; MLP]\label{eg:mlp}
    Define input, output, and parameter spaces as in \cref{eg:dln}.
    Let $\sigma : \R \to \R$ be an activation function.
    The parameter--function map sends $(A, B)$ to the function
        $f_{A,B} : \R^m \to \R^m$ such that for $x \in \R^m$,
    \begin{equation}
        f_{A,B}(x)
        = B \cdot \sigma( A x ),
    \end{equation}
    where we lift the activation function to operate element-wise over column
    vectors $A x \in \R^h$.
\end{example}
This kind of architecture is called a \textbf{multi-layer perceptron (MLP).}
Note that the two-layer DLN is recovered if we use the identity function as an
activation function. However, if we use a non-linear activation function, we
can index a much richer hypothesis class.

The expressivity of these architectures is limited, however, by the omission of
a basic detail---the inclusion of a bias parameter for each neuron. We invite
the reader to correct this omission as our first exercise, which serves as a
chance to familiarise oneself with the concept of a parameter--function map.

\begin{exercise}[Biased neurons]
\label{ex:biased-neurons}
    A biased neuron is a neuron with an additional parameter that is added to
    the weighted sum of its inputs before it produces its output. Bias
    parameters allow each layer to represent an \emph{affine} transform, rather
    than just a linear transform. For each of Examples \ref{eg:neuron},
    \ref{eg:linear}, \ref{eg:dln}, and \ref{eg:mlp}, extend the example to use
    biased neurons. Precisely define the parameter space and the
    parameter--function map in each case.
\end{exercise}

Modern deep learning leverages more involved parameter--function maps. Neural
network architectures are typically defined in a modular fashion, including
linear (or affine) layers and non-linear layers like the above, along with more
specialised layers (well-known examples including convolutional layers,
residual layers, and attention layers).

\subsection{Supervised deep learning and loss functions}

In a typical supervised deep learning setting, we are given a data set of pairs
    $(x^{(1)}, y^{(1)}),$ $\ldots,$ $(x^{(n)}, y^{(n)}) \in \X \times \Y$.
These pairs exemplify inputs and their corresponding (possibly noisy) outputs
from an unknown target function
    $f : \X \to \Y$
which we would like to approximate using a neural network. That is, we would
like to find a parameter $w \in \W$ such that the corresponding function $f_w$
approximates the target function $f$.

Formally, given an example $(x, y) \in \X \times \Y$, define a
\textbf{per-example loss} function
    $L_{x,y} : \W \to \R$
to measure the deviation of $f_w(x)$ from the output $y$. A typical loss
function for vector outputs $\Y \subset \R^m$ is to use the \textbf{squared
error} loss
\begin{equation}
    \label{eq:square-loss}
    L_{x,y}(w) = \|f_w(x) - y\|^2.
\end{equation}
If $\Y = \Delta^{m-1}$ is a space of discrete distributions over $m$ objects it
is typical to use \textbf{cross entropy} loss
\begin{equation}
    \label{eq:cross-entropy}
    L_{x,y}(w) = - \sum_{i=1}^m y(i) \log f_w(i \mid x).
\end{equation}

Given a per-example loss function and a \textbf{data distribution}
    $q \in \Delta(\X\times\Y)$
over the space of examples (representing our possibly-noisy target function),
we formalise the objective of supervised learning as finding $w \in \W$ so as
to minimise the \textbf{population loss} function
    $L : \W \to \R$
defined such that
\begin{equation}
    L(w) = \E_{(x,y) \sim q}[L_{x,y}(w)].
\end{equation}
The population loss aggregates differences on outputs for individual inputs
into an overall measure of difference between $f_w$ and the target function.

In practice, we often can't evaluate the population loss over the entire
input/output space. We instead optimise an estimator based on our data set of
$n$ input--output pairs assumed to be sampled independently and identically
from $q$. Define an \textbf{empirical loss} function
    $L_n : \W \to \R$
such that
\begin{equation}
    L_n(w) = \frac1n \sum_{i=1}^n L_{x^{(i)},y^{(i)}}(w).
\end{equation}
If the per-example loss is squared error, the empirical loss is known as the
\textbf{mean squared error} objective.
\begin{equation}
    \label{eq:mse}
    L_n(w)
    =
    \frac1n \sum_{i=1}^n \left\| f_w(x^{(i)}) - y^{(i)} \right\|^2.
\end{equation}

\begin{exercise}[Empirical loss and population loss]
    \label{ex:empirical-population-loss}
    Fix $w \in \W$. Prove the following properties of the relationship between
    the empirical loss $L_n(w)$ and the population loss $L(w)$.
\begin{enumerate}
    \item
        For any $n$, the empirical loss is an unbiased estimator of the
        population loss. That is,
        \begin{equation}
            \E_{(x^{(i)}, y^{(i)}) \sim q}[L_n(w)] - L(w)
            = 0.
        \end{equation}
    \item
        As $n \to \infty$, the empirical loss converges almost surely to the
        population loss.
\end{enumerate}
\end{exercise}

Given a loss function, a \textbf{training algorithm} is a search algorithm that
aims to find $w \in \W$ such that the loss is approximately minimised. Most
modern deep learning algorithms are variants of \textbf{stochastic gradient
descent (SGD)}, which implements an iterative gradient-based local search of
the parameter space using empirical loss on subsamples of the data set.

Through many impressive feats of computer science and hardware/software
engineering, we are able to run such training algorithms to find low-loss
parameters within parameter spaces with billions of dimensions. The details are
vitally important to the success of modern deep learning, but are beyond the
scope of this tutorial. We only mention SGD in order to emphasise that
\emph{any local search method depends intimately on the properties of the
parameter--function map.}



\subsection{Statistical models and parameter--distribution maps}
\label{sec:prelims:stats}

The previous sections introduce neural networks and supervised deep learning
within a framework of function approximation. An alternative approach is to
view supervised learning as statistical inference, and particularly Bayesian
inference. This statistical framework is commonly used within the SLT
literature, and we will encounter it within this tutorial.

The first step to upgrading our framework is to move from talking about
functions to talking about conditional distributions. For each neural network,
we should be able to determine not just the output corresponding to each
input, but a probability density measuring how likely each output is given each
input.

Formally, let $\D \subseteq \Delta(\Y)^\X$ represent a class of
conditional distributions over $\Y$.
A distributional neural network architecture comprises a
\textbf{parameter--distribution map}
    $\Psi : \W \to \D$.
It is sometimes convenient to refer to the conditional distribution
    $\Psi(w) : \X \to \Delta(\Y)$
using the notation
    $p_w : \X \to \Delta(\Y)$.
The probability mass/density of a given output $y \in \Y$ conditional on a
given input $x \in \X$ is typically denoted either
    $p_w(y \mid x)$ or $p(y \mid x, w)$.

Sometimes, defining a parameter--distribution map is simply a re-framing. For
example, in image classification, our output space $\Y$ was a space of
distributions over image classes. Let $\C$ represents the underlying class
space, then $\Y = \Delta(\C)$ and we can view our parameter--function map
$\Phi : \W \to \F$ as a parameter--distribution map $\Psi : \W \to \D$ with
    $\D = \F \subseteq \Y^\X = \Delta(\C)^\X$.

In other cases, we can easily construct a parameter--distribution map from a
parameter--function map by introducing a \textbf{noise model}. For example, in
regression with an output space
    $\Y = \R^m$,
given a parameter--function map $w \mapsto f_w$ we can define a
parameter--distribution map
    $\Psi : \W \to \Delta(\Y)^\X$
such that, for $w\in\W$ and $x\in\X$,
\begin{equation}
\label{eq:gaussian-model}
    p(y \mid x, w) = \frac{1}{(2\pi)^{m/2}} 
    \exp\!\left(-\frac{\|y - f_w(x)\|^2}{2}\right).
\end{equation}
The above \textbf{Gaussian noise model} amounts to modelling each output $y$ as
drawn given $x, w$ according to the rule $y = f_w(x) + \varepsilon$ with
$\varepsilon \sim \mathcal{N}(0, I_m)$.

As the following exercises explore, performing statistical inference according
to the principle of maximum likelihood aligns with the loss-minimisation
perspective (given an appropriate choice of loss function).

\begin{exercise}[Maximum likelihood as MSE minimisation]
    \label{ex:nll-mse}
    Consider a regression problem with $\X = \Y = \R^m$. Assume we have a
    neural network architecture with a parameter space $\W$. Form a
    corresponding statistical model using the Gaussian noise model of
    \eqref{eq:gaussian-model}.
    Let 
        $(x^{(1)}, y^{(1)}),$ $\ldots,$ $(x^{(n)}, y^{(n)}) \in \X \times \Y$
    be a data set of $n$ regression examples. 
\begin{enumerate}
    \item
        Write an expression for the conditional probability of
        observing the labels
            $Y_n = (y^{(1)}, \ldots, y^{(n)})$
        given the inputs
            $X_n = (x^{(1)}, \ldots, x^{(n)})$
        and some fixed parameter $w \in W$, assuming the labels are sampled
        independently.
    \item
        Denote the resulting quantity by
            $p(Y_n \mid X_n, w)$.
        It is called the \textbf{likelihood} of the data set according to the
        model $w$. Compute and simplify the \textbf{negative log-likelihood},
        given by $-\log p(Y_n \mid X_n, w)$.
    \item
        Show that if $L_n$ is the mean squared error loss function from
        \eqref{eq:mse}, then
        $$
            \argmax_{w \in \W} p(Y_n \mid X_n, w)
            =
            \argmin_{w \in \W} L_n(w)
        .$$
\end{enumerate}
\end{exercise}

\begin{exercise}[Maximum likelihood as cross entropy minimisation]
    \label{ex:nll-cross-entropy}
    Consider a classification problem with $\Y = \Delta^{m-1}$. Assume we have
    a neural network architecture with a parameter space $\W$. As discussed,
    the parameter--function map in this case corresponds to a
    parameter--distribution map to conditional distributions over the
    underlying discrete class space $\C = \{1, \ldots, m\}$.
    Let 
        $(x^{(1)}, c^{(1)}),$ $\ldots,$ $(x^{(n)}, c^{(n)}) \in \X \times \C$
    be a data set of $n$ labelled classification examples.
    For each $c^{(i)} \in \C$ define
    \begin{equation}\label{eq:classification-dirac}
        y^{(i)} = \delta_{c^{(i)}} \in \Delta^{m-1}
    \end{equation}
    to be the corresponding Dirac distribution (with
    probability mass $1$ concentrated on $c^{(i)}$).
\begin{enumerate}
    \item
        Write an expression for the likelihood of the data set,
            $p(C_n \mid X_n, w)$,
        that is, the conditional probability of observing the labels
            $C_n = (c^{(1)}, \ldots, c^{(n)})$
        given the inputs
            $X_n = (x^{(1)}, \ldots, x^{(n)})$
        and some fixed parameter $w \in W$. Assume the labels are sampled
        independently.
    \item
        Show that the negative log-likelihood
            $$
            -\log p(C_n \mid X_n, w)
            = - \sum_{i=1}^n \log f_w(c^{(i)} \mid x^{(i)})
            .
            $$
    \item
        Using \eqref{eq:classification-dirac}, show that if the per-example loss
        function for the optimisation problem is cross entropy loss
            \eqref{eq:cross-entropy},
        then
        $$
            \argmax_{w \in \W} p(C_n \mid X_n, w)
            =
            \argmin_{w \in \W} L_n(w)
        .$$
\end{enumerate}
\end{exercise}

\begin{exercise}[Maximum likelihood as loss minimisation, in general]
    \label{ex:nll-energy}
    Consider a neural network architecture with a parameter space $\W$,
    parameter--function map $\Phi$, and output space $\Y = \R^m$.
    Let $L_{x,y} : \W \to \R$ be a per-example loss function. Assume that for
    each $x \in \X$ and $w \in \W$, the function $y \mapsto \exp(-L_{x,y}(w))$
    is integrable over $\Y$. Define a corresponding statistical model with
    parameter--distribution map $\Psi$ mapping $w \in \W$ to
    $p_w : \X \to \Delta(\Y)$ such that for all $y \in \Y$, $x \in \X$, and
    $w \in \W$, we have
    $$
        p(y \mid x, w)
        =
        \frac
            {\exp\!\left(- L_{x,y}(w)\right)}
            {Z(x,w)}
        ,
        \qquad
        Z(x,w)
        =
        \int_{\Y}\exp\!\left(- L_{x,z}(w) \right)dz
        .
    $$
    We re-use the notation $X_n$, $Y_n$ from \cref{ex:nll-mse}, and assume
    that labels are sampled independently.
\begin{enumerate}
    \item
        Show that the negative log-likelihood decomposes as
        $$
            -\log p(Y_n \mid X_n, w)
            =
            \sum_{i=1}^n L_{x^{(i)},y^{(i)}}(w)
            +
            \sum_{i=1}^n \log Z(x^{(i)}, w)
        .$$
    \item
        Now suppose the loss function takes the form
            $L_{x,y}(w) = \ell(f_w(x) - y)$
        for some integrable function $\ell : \R^m \to \R$.
        Show that $Z(x, w)$ is independent of $w$, and conclude that
        $$
            \argmax_{w \in \W} p(Y_n \mid X_n, w)
            =
            \argmin_{w \in \W} L_n(w)
        .$$
\end{enumerate}
\end{exercise}
\begin{remark}[Negative log-likelihood loss]
    Examples~\ref{ex:nll-mse}, \ref{ex:nll-cross-entropy}, and
    \ref{ex:nll-energy} nominally show that some maximum likelihood estimation
    problems have the same sets of ideal solutions as some carefully-chosen
    optimisation problems ($\argmin = \argmax$).
    In general, that the global optima for two optimisation problems coincide
    is insufficient to show that practical optimisation methods will always
    find the same solutions: practical optimisation methods don't always find
    global optima.
    However, in the above cases, your derivation likely reveals a stronger
    result: the optimisation landscapes for the function approximation
    problem are related to the optimisation landscape for the maximum
    likelihood estimation problem by a strictly monotonically decreasing
    transform (an affine transform of a logarithm).
\end{remark}
        
\subsection{Bayesian deep learning}
\label{sec:bayesian-dl}

The principle of maximum likelihood, explored in the preceding exercises,
is one approach to statistical inference. An alternative approach is
\textbf{Bayesian inference}, which treats $w$ as a random variable and
maintains a probability distribution over $\W$ reflecting our uncertainty about
which parameter best explains the data. The majority of theoretical results
in SLT have been developed within this setting
    \citep[cf.,]{Watanabe2009,Watanabe2018}.

Bayesian learning begins with a \textbf{prior} distribution
    $\varphi \in \Delta(\W)$,
a probability density on $\W$ representing our beliefs about the parameter
$w$ before observing any data. Throughout this tutorial, we assume $\varphi$ is
positive and smooth on $\W$.

\begin{definition}[Posterior distribution, partition function, free energy]
\label{def:posterior}
\label{def:partition-function}
\label{def:free-energy}
    Given a data set of $n$ input--output pairs
        $(x^{(1)}, y^{(1)}),$ $\ldots,$ $(x^{(n)}, y^{(n)}) \in \X \times \Y$,
    Bayes' rule updates the prior into the \textbf{posterior distribution}
        $\pi_n \in \Delta(\W)$, a probability density on $\W$ defined by
    \begin{equation}
    \label{eq:posterior}
        \pi_n(w)
        = \frac{1}{Z_n} \varphi(w) \prod_{i=1}^{n} p(y^{(i)} \mid x^{(i)}, w)
    \end{equation}
    where $Z_n \in \R$ is a normalising constant called the \textbf{partition
    function} (or \textbf{marginal likelihood} or \textbf{model evidence}),
    \begin{equation}
    \label{eq:partition-function}
        Z_n
        = \int_{\W} \varphi(w) \prod_{i=1}^{n} p(y^{(i)} \mid x^{(i)}, w)\, dw.
    \end{equation}
    We often work with the \textbf{(Bayesian) free energy}, the negative log
    of the model evidence,
    \begin{equation}
    \label{eq:free-energy}
        F_n
        = -\log Z_n.
    \end{equation}
\end{definition}

The posterior re-weights the prior by the likelihood of the observed data:
parameters $w$ under which the data is more probable receive higher posterior
density, while parameters under which the data is improbable are
down-weighted. 

The partition function $Z_n$ is the marginal probability of observing the data
set, averaged over all parameters weighted by the prior. It quantifies how well
the statistical model \emph{as a whole} predicts the observed data. The free
energy is an inverted measure of how well the model fits the data.

\begin{exercise}[Bayesian posterior as a Gibbs distribution]
\label{ex:gibbs}
    To interpret Bayesian inference in terms more similar to function
    approximation, we introduce a loss function based on the
    \textbf{negative log-likelihood},
    \begin{align}
    \label{eq:nll-loss}
        L_{x,y}(w)
        &= -\log p(y \mid x, w),
    \\  L_n(w)
        &= -\frac{1}{n} \sum_{i=1}^{n} \log p(y^{(i)} \mid x^{(i)}, w).
    \end{align}
    Here, $p(y \mid x, w)$ is the conditional density from the
    parameter--distribution map and 
        $(x^{(1)}, y^{(1)}),$ $\ldots,$ $(x^{(n)}, y^{(n)}) \in \X \times \Y$
    is a data set of $n$ examples.
    Show that
    $$
        \pi_n(w)
        = \frac{1}{Z_n} \varphi(w) \exp\!\big({-}n L_n(w)\big)
    $$
    and
    $$
        Z_n
        = \int_{\W} \varphi(w) \exp\!\big({-}n L_n(w)\big)\, dw.
    $$
\end{exercise}

The \textbf{Bayesian deep learning process} is the process of re-weighting each
individual parameter vector in response to seeing an increasing amount of data,
producing the sequence of belief distributions
    $\varphi, \pi_1, \pi_2, \ldots \in \Delta(\W)$.
Over time, the posterior will concentrate around parameters which provide a good
fit for the data (though there is more to the story in the singular case, as we
will see in later sections).

Bayesian deep learning is a \emph{global search} in that all parameter vectors are
considered in parallel. This makes exact Bayesian learning computationally
intractable for non-trivial models, but more analytically tractable than SGD in
general. Moreover, as we will explore in later sections, the dynamics of Bayesian
learning still reflect the local geometry of the parameter--distribution map.

In preparation for exploring this connection between geometry and learning we
finally introduce localised variants of the partition function and free energy,
where the integral in (\ref{eq:partition-function}) is restricted to a certain
neighbourhood in parameter space.

\begin{definition}[Local partition function and local free energy]
\label{def:local-free-energy}
    Given a neighbourhood $\U \subseteq \W$, the
    \textbf{local partition function} and \textbf{local free energy} are
    \begin{equation}
    \label{eq:local-free-energy}
        Z_n(\U) = \int_{\U} \varphi(w) \prod_{i=1}^{n} p(y^{(i)} \mid x^{(i)}, w)\, dw,
        \qquad
        F_n(\U) = -\log Z_n(\U).
    \end{equation}
\end{definition}

The local partition function (free energy) measures how well (poorly) parameters
within a given neighbourhood collectively explain the data.
Note, $Z_n(\W) = Z_n$ and $F_n(\W) = F_n$.

\begin{exercise}[Local posterior mass and local free energy]
    \label{ex:posterior-free-energy}
    Let $\U, \V \subseteq \W$ be two neighbourhoods of parameter space.
    Let $\pi_n(\U) = \int_{\U} \pi_n(w) \,dw$ denote the posterior mass of
    a neighbourhood. Show the following.
\begin{enumerate}
    \item $\displaystyle \pi_n(\U) = \frac{Z_n(\U)}{Z_n}$ (the same holds for $\V$).
    \item $\displaystyle \log \frac{\pi_n(\U)}{\pi_n(\V)} = F_n(\V) - F_n(\U)$.
\end{enumerate}
In conclusion, the posterior odds ratio of the two regions
    $\pi_n(\U) / \pi_n(\V)$
depends exponentially on the difference between their local free energies.
\end{exercise}

\begin{exercise}[Partition function of partitioned parameter space]
    \label{ex:partition-decomposition}
    Let $\U_1, \ldots, \U_k \subseteq \W$ be disjoint subsets with
    $\W = \U_1 \sqcup \cdots \sqcup \U_k$.
\begin{enumerate}
    \item
        \label{ex:partition-decomposition:logsumexp}
        Show that the partition function decomposes as
        $Z_n = \sum_{j=1}^k Z_n(\U_j)$, and hence
        \begin{equation*}
            F_n
            = -\log \sum_{j=1}^k e^{-F_n(\U_j)}.
        \end{equation*}
    \item
        Show that if $F_n(\U_1) < F_n(\U_j)$ for all $j \neq 1$, then
        \begin{equation*}
            F_n(\U_1) - \log k < F_n < F_n(\U_1).
        \end{equation*}
\end{enumerate}
In conclusion, the overall free energy is dominated by the region(s) with the
lowest local free energy.
\end{exercise}

\clearpage
\section{What is degeneracy?}
\label{sec:degeneracy}

In this section, we define degeneracy as a property of parameter--function maps, and study the relationship between this property and other similar notions.

\subsection{A definition of degeneracy}
\label{sec:degeneracy:definition}

Consider a neural network architecture with parameter space $\W \subseteq \R^d$ and parameter--function map
    $\Phi : \W \to \F$ ($w \mapsto f_w$).
Say the parameter--function map is \textbf{degenerate at $w\in\W$} if there is a non-zero vector $v \in \R^d$ such that the directional derivative of the parameter--function map in the direction $v$ is zero:
\begin{equation}\label{eq:degeneracy}
    D_v \Phi(w) = 0.
\end{equation}

The directional derivative $D_v \Phi(w)$ for non-zero $v \in \R^d$ may be defined variously as
\begin{equation}\label{eq:directional-derivative}
    D_v \Phi(w)
    = \frac{v}{\|v\|} \cdot \nabla \Phi(w)
    = \sum_{i=1}^d \frac{v_i}{\|v\|} \frac{\partial \Phi}{\partial w_i}(w)
    = \lim_{\epsilon \to 0} \frac{\Phi(w + \epsilon v) - \Phi(w)}{\epsilon\|v\|}.
\end{equation}
We conventionally define $D_0 \Phi(w) = 0$.
Note that $w$ is fixed, but the directional derivative is still a function (of
the same type as $\Phi(w)$; $D_v \Phi(w) : \X \to \Y$), since it represents an
infinitesimal change in $\Phi(w)$ in response to a change in $w$ in the
direction $v$. In \eqref{eq:degeneracy} we mean that this function is
identically zero for all inputs $x \in \X$ for the given $w \in \W$.

Intuitively, the parameter--function map is degenerate at $w \in \W$ if there
is a direction in which we can infinitesimally perturb $w$ without changing the
function $f_w$.

This definition of degeneracy applies to a single point in parameter space. As we shall soon see, it may be the case that the parameter--function map is degenerate at some points but not at others. We can clarify the situation with new terminology.
\begin{itemize}
    \item We say that a parameter--function map is \textbf{somewhere degenerate} if it is degenerate at \emph{any} point in parameter space.
    \item We say that a parameter--function map is \textbf{everywhere degenerate} if it is degenerate at \emph{all} points in parameter space.
\end{itemize}

By convention, if we say that the parameter--function map is simply
\textbf{degenerate} (without specifying somewhere, everywhere, or at a
particular point), we mean that it is \emph{somewhere degenerate}.

The following exercises explore this definition in some toy parametrisations of a simple function class (constants), displaying in the simplest possible setting some basic forms of degeneracy that will arise repeatedly throughout the tutorial.

\begin{exercise}[Parametrising the space of constants]
\label{ex:univariate-degeneracy}
Let $\X = \{\ast\}$ and $\Y = \R$, so that we have a hypothesis class of constants. Consider the scalar parameter space $\W = \R$ and the parameter function map that maps $w$ to $f_w = w$ (the output is just the parameter itself).
\begin{enumerate}
    \item What is the directional derivative of the parameter--function map in
        direction $v = 1$?
    \begin{hint}
What kind of derivative does this reduce to?
    \end{hint}
    \item At which points in the parameter space is this parameter--function map degenerate, if any?
\end{enumerate}
\end{exercise}

\begin{exercise}[Degeneracy from raising a parameter to a power]
\label{ex:cubic-degeneracy}
Again, consider a hypothesis class of constants and the scalar parameter space $\W = \R$. This time, define a parameter function map that maps $w$ to $f_w = w^3$.
\begin{enumerate}
    \item Show that this architecture indexes exactly the same hypothesis class as the architecture described in \cref{ex:univariate-degeneracy}.
    \item What is the directional derivative of the parameter--function map in
        direction $v = 1$?
    \begin{hint}
What kind of derivative does this reduce to?
    \end{hint}
    \item At which points in the parameter space is the parameter--function map degenerate, if any?
\end{enumerate}
\end{exercise}

\begin{exercise}[Degeneracy from multiplying two parameters together]
\label{ex:product-degeneracy}
Again, consider a hypothesis class of constants. Consider this time the two-dimensional parameter space $\W = \R^2$. Define a parameter function map that maps $w = (a, b)$ to $f_w = a \cdot b$.
\begin{enumerate}
    \item
        Show that this architecture indexes exactly the same hypothesis class
        as the architecture described in \cref{ex:univariate-degeneracy}.
    \item
        What is the directional derivative of the parameter--function map in
        direction $v = (1, 0)$?
        \begin{hint}
what kind of derivative does this reduce to?
        \end{hint}
    \item
        What is the directional derivative of the parameter--function map in
        direction $v = (0, 1)$?
        \begin{hint}
what kind of derivative does this reduce to?
        \end{hint}
    \item \label{ex:product-degeneracy:axes-directions}
        At which points in the parameter space is the parameter--function map
        degenerate \emph{in these directions,} if any?
    \item
        At which points in the parameter space is the parameter--function map
        degenerate, if any?
\end{enumerate}
\end{exercise}

\begin{exercise}[Degeneracy from adding two parameters together]
\label{ex:sum-degeneracy}
Again, consider a hypothesis class of constants. Consider this time the two-dimensional parameter space $\W = \R^2$. Define a parameter function map that maps $w = (a, b)$ to $f_w = a + b$.
\begin{enumerate}
    \item Show that this architecture indexes exactly the same hypothesis class as the architecture described in \cref{ex:univariate-degeneracy}.
    \item What is the directional derivative of the parameter--function map in
        direction $v = (1, -1)$?
    \item At which points in the parameter space is the parameter--function map degenerate, if any?
\end{enumerate}
\end{exercise}

\subsection{Degeneracy and continuous symmetries}

A common source of degeneracy in neural network architectures is the
presence of continuous symmetries of the parameter--function map. A continuous
symmetry traces out a curve of functionally equivalent parameters. The tangent
to this curve is a degenerate direction. In this section, we will explore this
kind of symmetry and some examples from deep learning.

A \textbf{symmetry} of a parameter--function map $\Phi$ is a transformation of
parameter space that maps parameters while preserving the implemented function.
That is, a transformation
    $T : \W \to \W$
is a symmetry if for all $w \in \W$, we have
    $\Phi(w) = \Phi(T(w))$ ($f_w = f_{T(w)}$).

A \textbf{continuous symmetry} of a parameter--function map $\Phi$ is a family
of transformations
    $\{T_t : \W \to \W\}_{t \in \R}$
indexed by a parameter $t\in\R$, such that
\begin{enumerate}[label=(\roman*)]
    \item $T_t$ is a symmetry ($\Phi \circ T_t = \Phi$) for all $t \in \R$,
    \item $T_0$ is the identity transformation on $\W$, and
    \item The map $t \mapsto T_t(w)$ is differentiable for each $w \in \W$.
\end{enumerate}
The continuous symmetry is \textbf{trivial at $w$} if
    $\left.\frac{d}{dt}\right|_{t=0}  T_t(w) = 0$,
or else \textbf{non-trivial at $w$}.

\begin{exercise}[Continuous symmetry example]
    \label{ex:sum-symmetry}
    Recall the parameter--function map from \cref{ex:sum-degeneracy}, with $\W
    = \R^2$ and with $(a, b) \in \W$ mapping to the constant function
    $f_{a,b} = a + b$.
    Consider the family of transformations $T_t : \W \to \W$ for $t \in \R$
    such that
        $$T_t(a, b) = (a + t, b - t).$$
    \begin{enumerate}
        \item Describe the effect of this family of transformations on the
            parameter space.
        \item Show that this family of transformations is a continuous symmetry.
        \item Show that this continuous symmetry is non-trivial everywhere.
    \end{enumerate}
\end{exercise}

At each parameter, a non-trivial continuous symmetry traces out a curve of
functionally equivalent parameters. This indicates the presence of a direction
in parameter space in which the function does not change. The following
proposition formalises this connection between continuous symmetries and
degeneracy.

\begin{proposition}
\label{prop:symmetry-degeneracy}
    If a parameter--function map $\Phi$ admits a continuous symmetry $\{T_t\}$
    that is non-trivial at $w$, then $\Phi$ is degenerate at $w$.
\end{proposition}

\begin{proof}
    Define the curve in parameter space traced by the continuous symmetry,
    $\gamma : \R \to \W$ with $\gamma(t) = T_t(w)$.
    Since $T_t$ is a continuous symmetry, the derivative along the curve
    vanishes for all $t$:
    \begin{align*}
        D_{\gamma'(t)} \Phi(\gamma(t))
        &\propto \gamma'(t) \cdot \nabla \Phi(\gamma(t))
    \\  &= \frac{d}{dt} \Phi(\gamma(t))
            \tag{chain rule}
    \\  &= \frac{d}{dt} \Phi(w)
            \tag{symmetry, $\Phi \circ T_t = \Phi$}
    \\  &= 0.
            \tag{$\Phi(w)$ is constant in $t$}
    \end{align*}
    In particular, at $t=0$, we have
        $0 = D_{\gamma'(0)} \Phi(\gamma(0)) = D_{\gamma'(0)} \Phi(w)$
    since $\gamma(0) = w$ by identity. Since $T_t$ is non-trivial at $w$,
    $\gamma'(0)$ is non-zero, so $\gamma'(0)$ is a degenerate direction at $w$.
\end{proof}

\begin{corollary}
    If a parameter--function map $\Phi$ admits a continuous symmetry $\{T_t\}$
    that is non-trivial at every parameter $w \in \W$, then $\Phi$ is
    everywhere degenerate.
\end{corollary}

The following exercises exhibit some continuous symmetries in two neural
network architectures, a small ReLU MLP and a toy autoencoder. Similar
architectures are studied in more detail from an SLT perspective by
\citet{Carroll2021} and \citet{Chen+2023} respectively.

\begin{exercise}[ReLU scaling symmetry]
\label{ex:relu-scaling}
Consider a two-layer MLP (\cref{eg:mlp}) with $\sigma = \relu$, scalar inputs and outputs ($m = 1$), and a single hidden unit ($h = 1$). The parameter space is $\W = \R^2$ with parameters $(a, b)$, and the parameter--function map is $f_{a,b}(x) = b \cdot \relu(ax)$.
\begin{enumerate}
    \item
        Show that $\relu$ is \emph{positively homogeneous}: $\relu(\alpha z) =
        \alpha \, \relu(z)$ for all $\alpha > 0$ and $z \in \R$.
    \item
        Define $T_t(a, b) = (e^t a, \, e^{-t} b)$ for $t \in \R$. Show that
        $\{T_t\}$ is a continuous symmetry of this parameter--function map.
    \item
        For a fixed parameter $(a, b) \neq (0, 0)$, compute the degenerate
        direction arising from this symmetry at $(a, b)$.
    \item
        Show that the symmetry is trivial at the origin. Is this
        parameter--function map degenerate at the origin?
\end{enumerate}
\end{exercise}

\begin{exercise}[Rotation symmetry in a linear autoencoder]
\label{ex:rotation-symmetry}
    Consider a \emph{linear autoencoder} with bottleneck dimension $h$ and ambient dimension $m \geq h$. The parameter space encodes a single matrix $W \in \R^{h \times m}$ and the parameter--function map sends $W$ to the function $\Phi(W) = f_W : \R^m \to \R^m$ where
    \begin{equation*}
        f_W(x) = W^\top W x
    \end{equation*}
    for $x \in \R^m$. Here, $W$ serves simultaneously as encoder ($x \mapsto Wx
    \in \R^h$) and decoder ($z \mapsto W^\top z \in \R^m$).
\begin{enumerate}
    \item
        Show that for any orthogonal matrix $R \in \R^{h \times h}$ (i.e., any
        $h \times h$ matrix such that $R^\top R = I$), the map $T_R(W) = RW$
        is a symmetry.
    \item
        Specialise to $h = 2$. Using the rotation matrices
        \begin{equation*}
            R(\theta) = \begin{pmatrix}
                \cos\theta & -\sin\theta
            \\  \sin\theta & \cos\theta
            \end{pmatrix},
        \end{equation*}
        show that $T_\theta(W) = R(\theta) W$ defines a continuous symmetry.
    \item
        For a fixed parameter $W \in \R^{2 \times m}$, compute the degenerate
        direction that arises from this symmetry.
    \item
        For general bottleneck dimension $h$, how many independent continuous
        symmetries does the orthogonal group $O(h)$ contribute?
        \begin{hint}
What is the dimension of $O(h)$?
        \end{hint}
\end{enumerate}
\end{exercise}

\subsection{Localised degeneracy}

We have seen that a non-trivial continuous symmetry traces out curves of
functionally equivalent parameters throughout the entire parameter space, and
the tangent directions to these curves are degenerate directions at every
point. However, not all degenerate directions arise from globally continuous
symmetries. In fact, the more interesting case is when degeneracy affects only
some parameters and not others.


Within certain subsets of parameter space, many neural network architectures
(including those studied in
    \cref{ex:relu-scaling,ex:rotation-symmetry})
exhibit degenerate directions that cannot be defined in terms of
continuous transformations that are globally symmetries. In this section, we
explore localised degeneracies of this kind.

We begin with some examples of transformations that are only continuous
symmetries within certain subsets of parameter space.

\begin{exercise}[Localised symmetry example]
\label{ex:localised-symmetry}
    Recall the parameter--function map from \cref{ex:product-degeneracy},
    with $\W = \R^2$ and with $(a, b) \in \W$ mapping to the constant function
    $f_{a,b} = a \cdot b$.
    Consider the two families of transformations
    $T^{(a)}_t, T^{(b)}_t : \W \to \W$ for $t \in \R$ such that
        $$
        T^{(a)}_t(a, b) = (a + t, b),
        \qquad\qquad
        T^{(b)}_t(a, b) = (a, b + t)
        .
        $$
    \begin{enumerate}
        \item Describe the effect of each family of transformations on the
            parameter space.
        \item In which subset of parameter space does each family of
            transformations constitute a symmetry for all $t\in\R$?
        \item Compare your results to your answer to
            \cref{ex:product-degeneracy:axes-directions}.
    \end{enumerate}
\end{exercise}

The above example is technically a two-layer DLN with $m=h=1$. Let us now
explore localised degeneracy in a general two-layer DLN and then in a
non-trivial MLP.

\begin{exercise}[Degeneracy in the two-layer DLN]
\label{ex:dln-degeneracy}
    Consider the two-layer DLN from \cref{eg:dln}, with $h=m$. The parameter
    space encodes two matrices $A, B \in \R^{m \times m}$, and the
    parameter--function map sends $(A, B)$ to the function
        $\Phi(A,B) = f_{A,B} : \R^m \to \R^m$ such that $f_{A,B}(x) = BAx$
    for $x \in \R^m$.
\begin{enumerate}
    \item
        Let $(\delta\!A, \delta\!B)$ be a unit perturbation in parameter space
        (a unit vector in the underlying parameter space $\R^{2m^2}$, decoded
        into a pair of matrices).
        Show that the directional derivative of the parameter--function map
        at $(A, B)$ in direction $(\delta\!A, \delta\!B)$ is the linear map
        \begin{equation*}
            D_{(\delta\!A, \delta\!B)} \Phi(A, B)
            = B\,\delta\!A + \delta\!B\, A.
        \end{equation*}
        That is,
            $D_{(\delta\!A, \delta\!B)} \Phi(A, B) x
            = (B\,\delta\!A + \delta\!B\, A) x$.
        \begin{hint}
Use the limit definition, \eqref{eq:directional-derivative}.
        \end{hint}
    \item
        \label{ex:dln-degeneracy:null-rank}
        Show that at the zero parameter $(A, B) = (0, 0)$, every direction
        in parameter space is degenerate.
        Count the number of dimensions in the subspace of degenerate directions.
        
    \item
        \label{ex:dln-degeneracy:full-rank}
        Suppose both $A$ and $B$ are invertible. Show that
        $(\delta\!A, \delta\!B)$ is a degenerate direction if and only if
        $\delta\!B = -B\,\delta\!A\, A^{-1}$.
        Count the number of dimensions in the subspace of degenerate directions.
        
    \item
        \label{ex:dln-degeneracy:general-rank}
        Now consider the general case. Fix $A$ and $B$, and let $r_A = \rank(A)$ and $r_B = \rank(B)$. Show that the dimensionality of the space of degenerate directions is
        $$
            m^2 + (m-r_A)(m-r_B).
        $$
        \begin{hint}
The following is a guide to one possible approach.
        \begin{enumerate}
            \item
                Describe the image of the linear map $T_A$
                sending
                    $\delta\!B \in \R^{m^2}$
                to
                    $\delta\!B\, A \in \R^{m^2}$
                in terms of the row space of $A$.
                Show that the dimensionality of this image (the rank of the linear map $T_A$) is $m r_A$.
            \item 
                Similarly, describe the image of the linear map $T_B$
                sending
                    $\delta\!A \in \R^{m^2}$
                to
                    $B\, \delta\!A \in \R^{m^2}$
                in terms of the column space of $B$.
                Show that dimensionality of this image (the rank of the linear map $T_B$) is $m r_B$.
            \item 
                Describe the intersection of the images of $T_A, T_B$ in terms of the row/column spaces of $A, B$. Show that the dimensionality of this intersection is $r_A r_B$.
            \item
                Consider a third linear map, $T_D$,
                sending
                    $(\delta\!A, \delta\!B) \in \R^{2m^2}$
                to 
                    $B\,\delta\!A + \delta\!B\, A \in \R^{m^2}$.
                Describe the image of this linear map in terms of those of $T_A$ and $T_B$.
                Compute the rank of this linear map using Grassmann's identity.
            \item 
                What is the nullity of $T_D$? Why is this the same as the dimensionality of the space of degenerate directions?
        \end{enumerate}
    \end{hint}
\end{enumerate}
\end{exercise}

\begin{exercise}[Redundant units in an MLP]
\label{ex:mlp-degeneracy}
    Consider a single-hidden-layer MLP with scalar inputs and outputs, $h$
    hidden units with activation function $\sigma : \R \to \R$, and an output
    bias.
    The parameter space is $\W = \R^{3h+1}$ with parameters
    $w = (a_1, b_1, c_1, \ldots, a_h, b_h, c_h, d)$, and the
    parameter--function map is
    \begin{equation*}
        f_w(x) = d + \sum_{i=1}^{h} a_i \, \sigma(b_i x + c_i).
    \end{equation*}
    Here, for each hidden unit $i$, $a_i$ is the outgoing weight, $b_i$ is the
    incoming weight, and $c_i$ is the bias. The output bias is $d$.
\begin{enumerate}
    \item
        Suppose $a_i = 0$ (unit~$i$ has zero outgoing weight).
        Show that the directions $\delta b_i$ and $\delta c_i$ are degenerate
        at $w$.
    \item
        Suppose $b_i = 0$ (unit~$i$ has zero incoming weight).
        Show that the direction
            $(\delta a_i, \delta d) = (1, -\sigma(c_i))$
        (with all other components zero) is degenerate at $w$.
    \item
        Suppose $h \geq 2$ and $(b_i, c_i) = (b_j, c_j)$ for some $i \neq j$
        (units $i$ and $j$ have the same incoming weights and biases).
        Show that the direction $(\delta a_i, \delta a_j) = (1, -1)$
        (with all other components zero) is degenerate at $w$.
\end{enumerate}
\end{exercise}

\begin{remark}[Localised symmetries in neural networks]
    Each degenerate direction identified in the exercises above corresponds to
    a localised symmetry: a curve of functionally equivalent parameters that
    continuously extends only within a subset of parameter space.
    \begin{enumerate}
        \item
            In the DLN, the global change-of-basis symmetry
            $(A, B) \mapsto (G A, B G^{-1})$
            for invertible $G \in \R^{m\times m}$ accounts for $m^2$ degenerate
            directions.
            These directions are present even when $A$ and $B$ have full
            rank. However, as we showed in \cref{ex:dln-degeneracy}, when
            $A$ or $B$ drops rank, $(m - r_A)(m - r_B)$ additional degenerate
            directions open up.
            These extra directions are localised to the rank-deficient region
            of parameter space.
        \item
            In the MLP, at a parameter with $a_i = 0$, the path
            $t \mapsto (\ldots, a_i, b_i + t, c_i, \ldots)$
            traces equivalent parameters within $\{a_i = 0\}$ (since $f_w$
            does not depend on $b_i$ when $a_i = 0$), but this does not
            extend to a symmetry at nearby parameters with $a_i \neq 0$.
            Similarly, at a parameter with $(b_i, c_i) = (b_j, c_j)$, the
            path transferring weight between units $i$ and $j$ is a localised
            version of the sum symmetry from \cref{ex:sum-symmetry}.
    \end{enumerate}
\end{remark}

\begin{remark}[Rank of a neural network parameter]
\label{rem:rank}
    In both exercises above, the degree of degeneracy at a parameter is
    controlled by the amount of redundant capacity in the network.
    For the two-layer DLN, this is measured by the rank of the product $BA$:
    a rank-$r$ linear map can be implemented with a hidden dimension of $r$,
    leaving $m - r$ dimensions redundant.
    For a general single-hidden-layer MLP, this idea generalises in that
    we can define the rank of a parameter $w$ as the minimum number of
    hidden units needed to implement $f_w$
        \citep{FarrugiaRoberts2022,FarrugiaRoberts2024}.
    In the linear case, this ``neural network rank'' coincides with the
    matrix rank of $BA$.
    In both settings, lower rank corresponds to a higher number of
    degenerate directions.
\end{remark}

\begin{remark}[Measure zero degenerate sets]
    When a parameter--function map is degenerate somewhere but not everywhere, it
    is often the case that it is degenerate only within a measure-zero subset
    of parameter space.
    This means that sampling parameters uniformly at random
    results in a degenerate parameter with probability zero.
    However, this does not mean that the degeneracies can be dismissed. The
    learning process applies a non-random selection pressure and may select
    parameters from a measure-zero set. Moreover, the existence of degenerate
    parameters can have practical consequences for nearby non-degenerate
    parameters, and the collective neighbourhoods of all degenerate parameters
    comprise a non-measure-zero subset of parameter space.
\end{remark}

\subsection{Degeneracy and information singularities}

In the preceding sections, we studied degeneracy as a property of
parameter--function maps. In the statistical framework introduced in
\cref{sec:prelims:stats}, the fundamental object is instead a
parameter--distribution map $\Psi : \W \to \D$ ($w \mapsto p_w : \X \to
\Delta(\Y)$). In this section, we extend the definition of degeneracy to
parameter--distribution maps and connect it to a classical quantity from
statistics, the Fisher information matrix. We will show that under mild
regularity conditions on the statistical model, the Fisher information matrix
is singular at $w$ (has zero eigenvalues) if and only if the parameter--distribution
map is degenerate at $w$.

Given a parameter--distribution map $\Psi : \W \to \D$, say that $\Psi$ is
\textbf{degenerate at $w$ in direction $v$} if the conditional density does
not change to first order:
\begin{equation}\label{eq:dist-degeneracy}
    D_v \Psi(w) = 0.
\end{equation}
Here, $D_v \Psi(w)$ is not a conditional distribution but a signed difference
between conditional distributions, effectively a real function of $x$ and $y$
(in particular, it may include negative density changes). We mean that this
function is zero for all $x \in \X$ and all $y \in \Y$.

\Cref{eq:dist-degeneracy} extends the definition of degeneracy from
\cref{eq:degeneracy}. If $\Psi$ is constructed from a parameter--function map
$\Phi$ via a noise model (as in \cref{sec:prelims:stats}), then degeneracy of
$\Phi$ at $w$ in direction $v$ implies degeneracy of $\Psi$ at $w$ in direction
$v$, since $p(y \mid x, w)$ depends on $w$ only through $f_w(x)$.

To relate degeneracy to the Fisher information matrix, we introduce two
standard definitions.

\begin{definition}[Score function]
\label{def:score}
    Let $\Psi : \W \to \D$ be a parameter--distribution map with densities
    $p(y \mid x, w)$ that are positive and differentiable in $w$.
    The \textbf{score function} at $w$ is
    \begin{equation}\label{eq:score}
        s(x, y, w) = \nabla_w \log p(y \mid x, w) \in \R^d.
    \end{equation}
    The \textbf{directional score} in non-zero direction $v \in \R^d$ is
    \begin{equation}\label{eq:directional-score}
        s_v(x, y, w)
        = D_v \log p(y \mid x, w)
        = \frac{v}{\|v\|} \cdot \nabla_w \log p(y \mid x, w).
    \end{equation}
    The directional score measures how sensitive the log-density is to
    perturbations of $w$ in direction $v$.
\end{definition}

\begin{exercise}[Basic properties of the score function]
\label{ex:score-properties}
    Let $\Psi : \W \to \D$ be a parameter--distribution map with positive
    densities $p(y \mid x, w)$, differentiable in $w$. Assume that for
    each $x$ and $w$, the operations $\nabla_w$ and
    $\int_\Y \cdot\, dy$ may be exchanged.
\begin{enumerate}
    \item
        Show that the score function satisfies the following identity:
        for each $x \in \X$, $y \in \Y$, and $w \in \W$,
        \begin{equation}\label{eq:score-identity}
            s(x,y,w) = \frac{\nabla_w p(y \mid x,w)}{p(y \mid x,w)}.
        \end{equation}
    \item
        Show that the expected score is zero: for each $x \in \X$ and
        $w \in \W$,
        \begin{equation}
            \E_{y \sim p(y \mid x, w)}\bigl[s(x, y, w)\bigr] = 0.
        \end{equation}
        \begin{hint}
Differentiate the identity
            $\int_\Y p(y \mid x, w) \, dy = 1$.
        \end{hint}
    \item\label{ex:score-properties:b}
        Show that if $\Psi$ is degenerate at $w$ in direction $v$, then
        the directional score vanishes identically:
            $s_v(x, y, w) = 0$
        for all $x \in \X$ and $y \in \Y$.
    \item\label{ex:score-properties:c}
        Show the converse: if the directional score $s_v(x, y, w) = 0$
        vanishes for all $x \in \X$ and $y \in \Y$ at $w$, then $\Psi$ is
        degenerate at $w$ in direction $v$.
\end{enumerate}
\end{exercise}

\begin{definition}[Fisher information matrix]
\label{def:fim}
    Let $\Psi : \W \to \D$ be a parameter--distribution map with score
    function $s$.
    Let $q(x)$ denote the marginal distribution of inputs.
    The \textbf{Fisher information matrix}
        $I : \W \to \R^{d \times d}$
    is defined by
    \begin{equation}\label{eq:fim}
        I(w) = \E_{x \sim q(x)}\,
            \E_{y \sim p(y \mid x, w)}\!\bigl[
                s(x,y,w)\, s(x,y,w)^\top
            \bigr].
    \end{equation}
\end{definition}

The Fisher information matrix is symmetric and positive semidefinite at
every $w$ (being an expectation of positive semidefinite matrices
$s\, s^\top$).
However, the Fisher information matrix may not be symmetric positive \emph{definite,} that is, it may be singular.

The following exercise shows that singularity of the Fisher information matrix is equivalent to degeneracy of the parameter--distribution map under mild regularity conditions.

\begin{exercise}[Fisher information and degeneracy]
\label{ex:fim-degeneracy}
    Let $\Psi : \W \to \D$ be a parameter--distribution map with densities
    $p(y \mid x, w)$. Assume that $p(y \mid x, w) > 0$ for all
    $y \in \Y$, $x \in \X$, $w \in \W$, that $w \mapsto p(y \mid x, w)$
    is differentiable, and that $q(x) > 0$ for all $x \in \X$.
\begin{enumerate}
    \item \label{ex:fim-degeneracy:a}
        Show that for any unit vector $v \in \R^d$,
        \begin{equation}\label{eq:fim-quadratic}
            v^\top I(w)\, v
            = \E_{x \sim q(x)}\,
              \E_{y \sim p(y \mid x, w)}\!\bigl[
                s_v(x, y, w)^2
              \bigr].
        \end{equation}
    \item
        Using \cref{ex:score-properties:b}, show that if $D_v \Psi(w) = 0$
        for some non-zero $v$, then $I(w)$ is not positive definite.
        \begin{hint}
Check the definition of positive definite.
        \end{hint}
    \item
        Using \cref{ex:score-properties:c}, show that if $I(w)$ is not
        positive definite, then there exists a non-zero $v$ for which
        $D_v \Psi(w) = 0$.
        \begin{hint}
Check the definition of positive definite.
        \end{hint}
    \item
        Conclude that
        \begin{equation}
            \ker I(w) = \bigl\{v \in \R^d : D_v \Psi(w) = 0\bigr\}.
        \end{equation}
        That is, the null space of the Fisher information matrix is
        exactly the space of degenerate directions of the
        parameter--distribution map.
\end{enumerate}
\end{exercise}

\begin{remark}[Watanabe's strictly singular models]
\label{rem:strictly-singular}
    \citet{Watanabe2009} defines a statistical model as \emph{strictly
    singular} if either of two conditions hold:
    \begin{enumerate}
        \item
            The Fisher information matrix $I(w)$ is singular for some $w \in \W$.
        \item
            The parameter--distribution map is not one-to-one, that is, for some $w, w' \in \W$ such that $w \neq w'$, we have $\Psi(w) = \Psi(w')$.
    \end{enumerate}
    By \cref{ex:fim-degeneracy}, under mild regularity conditions, the first
    condition is equivalent to our notion of (somewhere) degeneracy. The second
    condition is called non-identifiability.

    \citet{Watanabe2007,Watanabe2009} observes that many of the results from
    classical statistics crucially assume among their regularity conditions the
    parameter--distribution map is identifiable and the Fisher information
    matrix is nonsingular. Whereas, in a statistical model based on a
    non-trivial neural network (or any other statistical model involving
    hierarchical structure), it is typical for the Fisher information matrix to
    include singularities.
\end{remark}

\subsection{Degeneracy and the loss landscape}
\label{sec:degeneracy:hessian}

In the preceding sections, we studied degeneracy as a property of
parameter--function maps and parameter--distribution maps. We now consider the
implications of degeneracy in these maps for the \emph{loss landscape} in which
deep learning algorithms operate.

Recall the definitions of loss functions from \cref{sec:prelims}. In what
follows, we assume that the per-example loss $L_{x,y}(w)$ depends on $w$ only
through $\Phi(w)$, and work with the population loss
    $L(w) = \E_{(x,y) \sim q}[L_{x,y}(w)]$.
For statistical models,
we use the expected negative log likelihood
    $L(w) = \E_{(x,y) \sim q}[-\log p(y \mid x, w)]$
as our loss function. In either case, we assume that $L$ is
twice-differentiable with respect to $w$.

Let us begin with some basic observations about the relationship between degeneracy
in parameter--function maps and directional derivatives in the loss landscape.

\begin{exercise}[Directional derivatives of the loss]
\label{ex:loss-basic}
Let $L : \W \to \R$ be a population loss function satisfying the assumptions described above.
\begin{enumerate}
    \item
        Show that if $\Phi$ is degenerate at $w$ in direction $v$, then the
        directional derivative of the loss vanishes in the same direction:
        $D_v L(w) = 0$.
    \item
        Suppose $\W = \R^d$. Show that if $d > 1$, then for \emph{any} any $w \in \W$, there exists a nonzero $v$
        such that $D_v L(w) = 0$. \begin{hint}
consider separately the cases $\nabla
        L(w) = 0$ and $\nabla L(w) \neq 0$.
        \end{hint}
\end{enumerate}
\end{exercise}

We see that parameter--function map degeneracy implies flat loss directions,
but a vanishing directional derivative of the loss landscape is commonplace. To
find a satisfying definition of degeneracy in a loss landscape, we must look at
second-order information.

Define the \textbf{Hessian matrix} $H(w) \in \R^{d \times d}$ of the loss at
$w$ by
\begin{equation}
    H(w)_{jk} = \frac{\partial^2 L}{\partial w_j \partial w_k}(w).
\end{equation}
The Hessian at $w$ is sometimes denoted $\nabla^2 L(w)$.
The Hessian captures the local curvature of the loss landscape via the
quadratic approximation
\begin{equation}
    L(w + \delta)
    \approx L(w)
    + \nabla L(w)^\top \delta
    + \tfrac{1}{2}\,\delta^\top H(w)\, \delta.
\end{equation}
At a critical point ($\nabla L(w) = 0$), the Hessian determines whether the
loss curves upward, downward, or remains flat in each direction.

In particular, at a local minimum, $H(w)$ is positive semidefinite. We can
therefore classify local minima as follows:
\begin{itemize}
    \item A local minimum $w$ is \textbf{regular} (or \textbf{non-degenerate},
        or \textbf{Morse}) if $H(w)$ is positive definite.
    \item A local minimum $w$ is \textbf{degenerate} (or \textbf{non-Morse}) if
        $H(w)$ is singular (has a zero eigenvalue).
\end{itemize}

\begin{exercise}[Some examples of loss landscape degeneracy]
\label{ex:loss-landscape-degeneracy}
    Consider the parameter space $\W = \R^2$ with the identity
    parameter--function map $\Phi = \mathrm{id}$, so that we identify
    parameters with the functions they implement
    (cf., \cref{ex:univariate-degeneracy}).
    Consider the family of loss functions
        $L_{k,l}(a,b) = a^{2k} + b^{2l}$
    for non-negative integers $k$ and $l$.
\begin{enumerate}
    \item
        Show that the origin is a global minimum of $L_{k,l}$ for all
        non-negative $k$ and $l$. For which $k$ and $l$ is it the unique global
        minimum?
    \item
        \label{ex:lld:1}
        Show that if $k = l = 1$, then the origin is a non-degenerate
        (regular) minimum.
        \begin{hint}
Compute the Hessian.
        \end{hint}
    \item
        \label{ex:lld:2}
        Show that if $k > 1$ and $l \geq 1$, then the origin is a
        degenerate minimum.
    \item
        \label{ex:lld:3}
        Show that if $k = 0$ and $l \geq 1$, then the origin is a
        degenerate minimum.
\end{enumerate}
\end{exercise}

Note that the Hessian degeneracy in \cref{ex:lld:2,ex:lld:3} arise for
qualitatively different reasons. In \cref{ex:lld:3}, the loss is constant along
the $a$-direction to all orders---a genuinely flat direction. In
\cref{ex:lld:2}, the loss does increase along the $a$-direction, just more
slowly than quadratically (as $a^4$, $a^6$, etc., depending on $k$). Moreover,
different values of $k$ give different rates of increase. The Hessian cannot
distinguish any of these cases: flat, quartic, and so on all appear the same
from the perspective of the Hessian rank. We will later see approaches that can
distinguish these degrees of degeneracy.

So much for defining loss landscape degeneracy. What is the relationship
between this kind of degeneracy and degeneracy of the
parameter--function/distribution map?

The relationship is subtle, and generally depends on the choice of loss
function. One natural setting in which to study this relationship is the
\emph{realisable statistical model:} given a parameter--distribution map and a
parameter $w_0$, if we consider $p_{w_0}$ as the true data-generating
distribution, the population negative log-likelihood loss will have a global
minimum at $w_0$. The following exercise shows that in this setting, under mild
regularity assumptions on the statistical model, degeneracy in the
parameter--distribution map at $w_0$ implies that the global minimum $w_0$ is a
degenerate global minimum of the loss landscape.

\begin{exercise}[Realisable models and Hessian degeneracy]
\label{ex:fim-hessian}
    Let $\Psi : \W \to \D$ be a parameter--distribution map with positive
    densities $p(y \mid x, w)$, twice-differentiable in $w$. Suppose data is
    generated from a fixed true parameter $w_0 \in \W$, meaning
    $q(y \mid x) = p(y \mid x, w_0)$. Consider the population negative
    log-likelihood loss
    \begin{equation}
        L(w)
        = -\E_{x \sim q(x)}\,
            \E_{y \sim p(y \mid x, w_0)}\!\bigl[\log p(y \mid x, w)\bigr].
    \end{equation}
    Assume that for each $x$, the operations $\nabla_w$ (and $\nabla_w^2$)
    and $\int_\Y \cdot\, dy$ may be exchanged.
\begin{enumerate}
    \item
        \emph{(Bartlett identity.)}
        Show that for each $x \in \X$ and $w \in \W$,
        \begin{equation}\label{eq:bartlett}
            -\E_{y \sim p(y \mid x, w)}\!\bigl[
                \nabla_w^2 \log p(y \mid x, w)
            \bigr]
            =
            \E_{y \sim p(y \mid x, w)}\!\bigl[
                s(x,y,w)\, s(x,y,w)^\top
            \bigr],
        \end{equation}
        where $s$ is the score function from \cref{def:score}.
        \begin{hint}
differentiate the identity
            $\E_{y \sim p(y \mid x, w)}[s(x, y, w)] = 0$
            (established in \cref{ex:score-properties}) with respect to $w$.
        \end{hint}
    \item
        Show that the Hessian of $L$ at the true parameter $w_0$ equals the
        Fisher information matrix (\cref{def:fim}):
        \begin{equation}\label{eq:fim-hessian}
            H(w_0) = I(w_0).
        \end{equation}
        \begin{hint}
compute $H(w)
            = -\E_{x}\, \E_{y \sim p(y \mid x, w_0)}
                [\nabla_w^2 \log p(y \mid x, w)]$,
            evaluate at $w = w_0$, and apply part~(a).
        \end{hint}
    \item
        Using the result of \cref{ex:fim-degeneracy}, conclude: if $\Psi$ is
        degenerate at $w_0$ in direction $v$, then $H(w_0)\, v = 0$.
        Contrapositively, if $H(w_0)$ is positive definite, then $\Psi$ is
        non-degenerate at $w_0$.
\end{enumerate}
\end{exercise}

\Cref{ex:fim-hessian} shows that for the population negative log-likelihood at
the true parameter, degeneracy of the parameter--distribution map implies a
degenerate loss landscape. The converse does not hold in general, nor does the
forward direction hold at arbitrary points in parameter space. The following
exercise explores these subtleties through examples.

\begin{exercise}[Contrasting parametric versus loss degeneracy]
\label{ex:loss-degeneracy-examples}
\begin{enumerate}
    \item
        Observe that the identity parameter--function map on $\W = \R^2$
        is non-degenerate everywhere. With this in mind, what do your answers
        to \cref{ex:lld:1,ex:lld:2,ex:lld:3} exemplify about the logical
        relationship between parameter--function map degeneracy and loss
        landscape degeneracy?
    \item
        Consider again the identity parameter--function map on $\W = \R^2$.
        Consider the loss function $L_\circ(a,b) = (a^2 + b^2 - 1)^2$.
        Find a global minimum of $L_\circ$ and show that it is a degenerate
        critical point.
        What does this example reveal about the logical relationship between
        parameter--function map degeneracy and loss landscape degeneracy?
    \item
        Now consider the cubic parameter--function map
            $\Phi(\alpha, \beta) = (\alpha^3, \beta^3)$.
        \begin{enumerate}
            \item
                Argue that $\Phi$ is degenerate at the origin (cf.,
                \cref{ex:cubic-degeneracy}).
            \item
                If we take $L_{k,l}$ from \cref{ex:loss-landscape-degeneracy} to be
                defined on the constant function $(a, b)$ (rather than equivalent
                underlying parameter), then in this parameterisation, that loss
                function becomes
                    $L_{k,l}(\alpha,\beta) = (\alpha^3)^{2k} + (\beta^3)^{2l}$.
                Show that the origin is a degenerate minimum of $L_{k,l}$ for all
                $k \geq 0$ and $l \geq 0$ under this parametrisation.
            \item
                What does this example reveal about the logical relationship
                between parameter--function map degeneracy and loss landscape
                degeneracy?
        \end{enumerate}
    \item
        Consider the product parameter--function map
            $\Phi(a,b) = a \cdot b$
        from \cref{ex:product-degeneracy}, and define
        $L(a,b) = \frac{1}{2}(ab - 1)^2$.
        \begin{enumerate}
            \item
                Show that $(0,0)$ is a critical point and that $\Phi$ is
                degenerate in every direction at $(0,0)$.
            \item
                Compute the Hessian at $(0,0)$ and show it is nonsingular.
            \item
                What does this example reveal about the logical relationship
                between parameter--function map degeneracy and loss landscape
                degeneracy? Why does this not contradict \cref{ex:fim-hessian}?
        \end{enumerate}
\end{enumerate}
\end{exercise}


\clearpage
\section{The degeneracy hierarchy}
\label{sec:llc}

We have so far explored different kinds of qualitative degeneracy in the loss landscape. A natural question arises: \emph{how can we quantify the degree of degeneracy at a particular point in parameter space?} In this section, we introduce the local learning coefficient, a rich mathematical object that reflects the degree of degeneracy at a parameter, and we explore it from several perspectives.

In this section and \cref{sec:fef} we assume that any population loss $L$ is a real analytic function. The precise definition of real analyticity is not important for this tutorial but interested readers can refer to \cite{Watanabe2009} (section 2). A property that will be useful is that such functions have a convergent Taylor expansion, in particular they have smooth derivatives of all orders.


\subsection{The local learning coefficient via volume scaling asymptotics}
\label{sec:llc:volume-scaling}

The key idea is to measure the \emph{volume} of near-optimal parameters. Consider a local minimum $w^* \in \W$ of the population loss $L$, and suppose we have a sufficiently small closed ball $B(w^*)$ centred on $w^*$ such that $L(w) \geq L(w^*)$ for all $w \in B(w^*)$. Given a tolerance $\epsilon > 0$, define the \textbf{sublevel set}
\begin{equation}\label{eq:sublevel-set}
    B(w^*, \epsilon) = \{ w \in B(w^*) \mid L(w) - L(w^*) < \epsilon \},
\end{equation}
consisting of all parameters in the ball whose loss is within $\epsilon$ of the minimum. The volume of this set is
\begin{equation}
    V(\epsilon) := \mathrm{Vol}(B(w^*, \epsilon)) = \int_{B(w^*, \epsilon)} dw.
\end{equation}

We will look at how this volume scales as $\epsilon \to 0$. It will be instructive for us to first calculate this scaling in the non-degenerate case.

\begin{exercise}[Volume scaling in the regular case]
\label{ex:llc-regular-case}
    Let $w^* \in \W$ be a local minimum of $L$, and suppose that the Hessian $\nabla^2 L(w^*)$ is positive definite. Show that
    \begin{equation} \label{eq:regular-volume-scaling}
        V(\epsilon) = c \, \epsilon^{d/2} + o(\epsilon^{d/2}) \quad \text{as } \epsilon \to 0,
    \end{equation}
    for some constant $c>0$ where $d = \dim(\W)$. You may assume that, as $\epsilon \to 0$, the volume of $B(w^*,\epsilon)$ agrees to leading order with the volume of the quadratic approximation obtained by replacing $L$ with its second-order Taylor expansion at $w^*$. 

    \begin{hint}
The volume of the ellipsoid $\{\boldsymbol{x} \in \R^d \mid \boldsymbol{x}^\top A \boldsymbol{x} < 1\}$ is proportional to $(\det A)^{-1/2}$.
    \end{hint}
\end{exercise}

The key takeaway from this exercise is that the dimension of the model controls the leading order asymptotics of the volume scaling. We can investigate this further by calculating the volume scaling in an example where the Hessian is not positive definite.

\begin{exercise}[Quadratic Valley Volume Scaling]
    \label{ex:quad-valley-scaling}
    Let $\W = [-R,R]^d$ for some (large) $R\in\R$ and let $L(x_1,\dots,x_d) = \sum\limits_{i=1}^{d'}x_i^2$.
    \begin{enumerate}
        \item Sketch the loss landscape when $d=2, \,d'=1$.
        \item $L$ achieves its minimum value $0$ at multiple points in $\W$. What is the dimension of the space $\W_0 :=  \{w\in\W \, | \,L(w)=0\}$? 
        \item Calculate the Hessian of $L$.
        \item Prove that $V(\epsilon) \propto \epsilon^{d'/2}$
        
        \begin{hint}
Use the volume of the ellipsoid from Exercise \ref{ex:llc-regular-case}
        \end{hint}
    \end{enumerate}
\end{exercise}

In this example, the degenerate directions do not contribute to the leading exponent of $V(\epsilon)$. We see that the exponent $\frac{d}{2}$ from Exercise \ref{ex:llc-regular-case} is replaced with $\frac{d'}{2}$ where $d'$ is the codimension of $\W_0$ i.e. the \textit{effective dimension} of $L$ which ignores the degenerate directions. This suggests that we can capture information about the degree of degeneracy of a loss function by looking at how the volume $V(\epsilon)$ scales as $\epsilon \to 0$. \cite{Watanabe2009} proves a general form for this kind of volume scaling which has been adapted by \cite{Lau+2025} into the following definition.

\begin{definition}[Local learning coefficient]
\label{def:llc}
    Let $w^* \in \W$ be a local minimum of the population loss $L$. There exists a unique rational number $\lambda(w^*) > 0$, a positive integer $\mu(w^*)$, and a constant $c > 0$ such that as $\epsilon \to 0$,
    \begin{equation}
    \label{eq:llc-volume}
        V(\epsilon) = c \, \epsilon^{\lambda(w^*)} (-\log \epsilon)^{\mu(w^*)-1} + o\!\left(\epsilon^{\lambda(w^*)} (-\log \epsilon)^{\mu(w^*)-1}\right).
    \end{equation}
    We call $\lambda(w^*)$ the \textbf{local learning coefficient (LLC)} at $w^*$, and $\mu(w^*)$ the \textbf{local multiplicity}.
\end{definition}

\begin{exercise}[Two equations for $\lambda$] \label{ex:solving-for-lambda}
\begin{enumerate}
    \item Prove that for $a>1$
    \begin{equation} \label{eq:llc-formula}
    \lambda = \lim_{\epsilon \to 0} \frac{\log(V(a\epsilon)/V(\epsilon))}{\log(a)}
    \end{equation}
    
    \item Prove that \begin{equation} \label{eq:solve-LLC}
    \lambda = \lim\limits_{\epsilon\to 0} \frac{\log V(\epsilon)}{\log \epsilon} \end{equation}
\end{enumerate}
\end{exercise}

Note that from equation \eqref{eq:regular-volume-scaling} we already have that the LLC and local multiplicity in the regular case are $\lambda = \frac{d}{2}$, $\mu=1$

When $\mu(w^*) = 1$, the formula simplifies to $V(\epsilon) = c\epsilon^{\lambda(w^*)} + o(\epsilon^{\lambda(w^*)})$ which allows us to interpret $\lambda(w^*)$ as the \textit{volume scaling exponent}: increasing the error tolerance by a factor of $a$ increases the volume of near-optimal parameters by a factor of $a^{\lambda(w^*)}$.

\begin{exercise}[Calculating the LLC]
\label{ex:calculating-llc}
    Calculate the LLC and the rank of the Hessian at the global minimum in the following cases:
    \begin{enumerate}
        \item $L(x)=x^2$
        \item $L(x) = x^{2m}$
        \item $L(x,y) = x^4 + y^4$ 
        \item \label{ex:calculating-LLC-d} $L(x,y) = x^2+y^4$
    \end{enumerate}
\end{exercise}

\begin{exercise}[Upper bound on the LLC]
\label{ex:lambda-upper-bound}
    In \cref{ex:llc-regular-case}, we saw that $\lambda = d/2$ for regular models. In this exercise, you will prove directly from the volume scaling definition that for any local minimum $w^*$ of the population loss, the local learning coefficient satisfies $\lambda(w^*) \leq d/2$.

    \begin{enumerate}
        \item 
             Let $\Lambda>0$ be an upper bound on the eigenvalues of the Hessian $\nabla^2 L(w)$ for all $w \in B(w^*)$. Using the Lagrange remainder form of Taylor's theorem, show that for all $w \in B(w^*)$,
            \begin{equation*}
                L(w) - L(w^*) \leq \frac{1}{2} \Lambda \|w - w^*\|^2.
            \end{equation*}
            \begin{hint}
for any real symmetric matrix $H\in\R^{n\times n}$ and vector $v\in \R^n$, $v^\top H v \leq \Lambda  \|v\|^2$ where $\Lambda$ is the largest eigenvalue of $H$.
            \end{hint}
            
        \item 
            Consider the standard Euclidean ball of radius $r$ centred at $w^*$, denoted $B_r(w^*) = \{w \in \W \mid \|w - w^*\| < r\}$. Find a radius $r$ (as a function of $\epsilon$ and $\Lambda$) such that for sufficiently small $\epsilon$,
            \begin{equation*}
                B_r(w^*) \subseteq B(w^*, \epsilon).
            \end{equation*}
            
        \item 
            Recall that the volume of a $d$-dimensional Euclidean ball of radius $r$ is proportional to $r^d$. Use your result from part (b) to show that there exists a constant $C > 0$ such that for sufficiently small $\epsilon$,
            \begin{equation*}
                  C \epsilon^{d/2} \leq V(\epsilon).
            \end{equation*}
            
        \item 
            Conclude that $\lambda(w^*) \leq d/2$. 
    \end{enumerate} 
    
\end{exercise}

We have now shown that the LLC satisfies $\lambda(w^*) \leq d/2$, with equality precisely in the regular (non-degenerate) case.
 It is also possible (although quite fiddly) to prove that $ \frac{r}{2} \leq \lambda(w^*)$ where $r= \rank \nabla^2L(w^*)$ and as we will see in the next exercise, the LLC captures more information about the geometry than just the Hessian rank.

\begin{exercise}[LLC vs Hessian Rank]
\label{ex:LLC-hessian-rank}
    In this exercise we will show that the LLC can detect differences in the local geometry that is not reflected in the rank of the Hessian.  
    
    Construct two loss functions $L_1$ and $L_2$ with minima $w_1$ and $w_2$ respectively such that $\rank \, \nabla^2L_1(w_1) = \rank \, \nabla^2L_2(w_2)$ but $\lambda(w_1) \neq \lambda(w_2)$.
\end{exercise}


\begin{exercise}[The cubically-parameterised loss]
\label{ex:cubic-parameterisation}
    We can make a quadratic loss degenerate by changing the parameterisation. Let the \emph{cubically-parameterised loss} be
    \begin{equation}
        L(\mu) = \tfrac{1}{2}(\mu^3 - \mu_0^3)^2, \qquad \mu \in \R
    \end{equation}
    obtained from the ordinary quadratic loss $\ell(\theta) = \frac{1}{2}(\theta - \theta_0)^2$ via the reparameterisation $\theta = \mu^3$, $\theta_0 = \mu_0^3$. 
    
    \begin{enumerate}
        \item
            Show that the cubically-parameterised loss is just as expressive as the ordinary quadratic loss: that is, they both have the same set of achievable minima.
        \item
            Compute the Hessian $L''(\mu_0)$ to show that the loss is degenerate at its minimum when $\mu_0 = 0$, and non-degenerate when $\mu_0 \neq 0$.
        \item
            Give an explicit formula for $V(\epsilon)$ when $\mu_0 = 0$ and $\mu_0 \neq 0$ separately.
        \item
            \label{ex:cubic-parameterisation:d}
            Use your answer to part~(c) to find the learning coefficient for $\mu_0=0$ and for $\mu_0\neq0$. Given that the model is non-degenerate when $\mu_0 \neq 0$, we expect the learning coefficient to be $d/2$ in that case; compare your answer. How does the cubically-parameterised loss at $\mu_0 = 0$ differ from the ordinary quadratic loss?
        \item
            \label{ex:cubic-parameterisation:e}
            Instead of taking $\epsilon \to 0$ to get the learning coefficient, fix a small but nonzero value for $\epsilon$, such as $\epsilon = 0.01$. Plot $V(\epsilon)$ as a function of $\mu_0$. As we saw from (d), the learning coefficient changes discontinuously when $\mu_0 = 0$---what happens with $V(\epsilon)$ as $\mu_0$ gets close to zero? What changes if you make $\epsilon$ smaller or larger?

        Even though the asymptotic \emph{learning coefficient} ($\epsilon \to 0$) only changes when $\mu_0 = 0$ exactly, note how the non-asymptotic \emph{volume} ($\epsilon$ finite) is affected in a larger neighbourhood.
    \end{enumerate}
\end{exercise}

\subsection{Perspectives on the local learning coefficient}

In this section, we survey several alternative perspectives on the LLC, from information theory, algebraic geometry, and geometric measure theory.

\paragraph{Information-theoretic interpretation.}

The LLC also admits a natural interpretation in terms of information \citep{Lau+2025}. The number of bits needed to specify the sublevel set $B(w^*, \epsilon)$ within the ball $B(w^*)$ is
\begin{equation}
    -\log_2 \frac{V(\epsilon)}{\mathrm{Vol}(B(w^*))}.
\end{equation}

For small $\epsilon$ and $\mu(w^*) = 1$, this is approximately $-\lambda(w^*) \log_2 \epsilon + O(\log_2 \log_2 \epsilon)$. In particular, the number of \emph{additional} bits needed to halve an already small error tolerance from $\epsilon$ to $\epsilon/2$ is
\begin{equation}
    -\log_2 \frac{V(\epsilon / 2)}{V(\epsilon)} \approx \lambda(w^*).
\end{equation}

That is, the LLC $\lambda(w^*)$ measures the number of bits required to halve the error near $w^*$. 

\paragraph{Local learning coefficient via algebraic geometry}
\label{par:rlct}

We explore an algebro-geometric perspective on the learning coefficient due to \citet{Watanabe2009}.
Note that this section is more advanced than the remainder of this tutorial, and should be considered optional.
When $L(w)$ is a real analytic function, Hironaka's resolution of
singularities theorem guarantees the existence of a smooth map
$g \colon M \to \mathcal{W}$ from an analytic manifold $M$ and local coordinates
$u = (u_1, \ldots, u_d)$ that \emph{monomialise} the loss near $w^*$:
\begin{equation*}
    L(g(u)) - L(w^*) = u_1^{2k_1} \cdots u_d^{2k_d},
    \qquad
    dw = b(u)\, u_1^{h_1} \cdots u_d^{h_d}\, du,
\end{equation*}
where $b(u) > 0$ is smooth and the exponents $k_j, h_j$ are
non-negative integers.
In these coordinates, the volume integral $V(\epsilon)$ can be
evaluated explicitly, yielding the asymptotic form of
\cref{def:llc} and giving
\begin{equation}\label{eq:rlct}
    \lambda(w^*) \;=\; \min_{\text{coordinate charts}}\; \min_{j:\, k_j > 0}
    \frac{h_j + 1}{2k_j}\,.
\end{equation}

The resolution map $g$ is far from unique, and different choices
produce different exponents $k_j, h_j$.
$\lambda(w^*)$ is nevertheless well-defined which can be seen from (another) equivalent definition using the \textbf{local zeta function}:
\begin{equation}\label{eq:local-zeta}
    \zeta(z,\, w^*) \;=\; \int_{B(w^*)}
        \bigl(L(w) - L(w^*)\bigr)^{\!z} \, dw\,,
\end{equation}
which converges for $\operatorname{Re}(z) > 0$ and extends to a
meromorphic function on $\mathbb{C}$ whose poles are all real,
negative, and rational.
The LLC $\lambda(w^*)$ is precisely the negative of the largest pole, and
the multiplicity $\mu(w^*)$ is the multiplicity of this pole.
Since \eqref{eq:local-zeta} makes no reference to any resolution,
$\lambda(w^*)$ depends only on $L$ and $w^*$. 

In the algebraic geometry literature, $\lambda(w^*)$ is called the
\textbf{real log canonical threshold (RLCT)} of $L$ at
$w^*$.
We collect several consequences of this perspective.

\begin{enumerate}
    \item The fact that the LLC is a rational number follows directly from equation \eqref{eq:rlct}.
    \item Equation \eqref{eq:rlct} suggests that we can compute 
        $\lambda(w^*)$ analytically by constructing an explicit resolution
        of singularities for $L(w) - L(w^*)$ and reading off the
        exponents.
        This is difficult in practice for neural network loss
        functions, but has been achieved for some architectures, as
        we will see in \cref{sec:llc:dln}.

    \item
        In \cref{sec:fef} we will state Watanabe's free energy formula
        connecting the Bayesian free energy to the LLC\@. The proof of this
        result makes repeated use of the monomialisation of the loss.
\end{enumerate}

\begin{exercise}[Computing the LLC via the RLCT]
\label{ex:rlct-computation}
    Suppose that $L(w) = w_1^{2k_1} \cdots w_d^{2k_d}$ near
    $w^* = 0$, where $k_j \geq 0$ are integers with at least
    one $k_j > 0$. Show that
    \[
        \lambda(0) = \min_{j:\, k_j > 0} \frac{1}{2k_j}\,.
    \]
\end{exercise}

\paragraph{Local learning coefficient as a fractal dimension} 

The LLC admits a geometric interpretation: it is a \emph{fractal dimension} of the parameter space, as measured through the lens of the loss function. To make this precise, we recall a standard notion from fractal analysis.

\begin{definition}[Hölder exponent]
\label{def:holder-exponent}
    Let $(X, \rho)$ be a (pseudo-)metric space and $\nu$ a measure on $X$. The \textbf{Hölder exponent} (or \textbf{local dimension}) $\alpha(x)$ of $\nu$ at a point $x \in X$ is given by
    \begin{equation} \label{eq:holder-exp}
        \alpha(x) = \lim_{\epsilon \to 0} \frac{\log \nu(B_\epsilon(x))}{\log \epsilon},
    \end{equation}
    where $B_\epsilon(x)$ is the ball of radius $\epsilon$ centred at $x$. Equivalently, $\alpha(x)$ is the unique real number such that 
    \begin{equation*}
        \nu(B_\epsilon(x)) = O(\epsilon^{\alpha(x)}) \qquad \text{as } \epsilon \to 0.
    \end{equation*}
\end{definition}

The Hölder exponent generalises the notion of dimension: when $\nu$ is the Lebesgue measure on $\R^d$ and $\rho$ is the Euclidean metric, every point has $\alpha(x) = d$, but for measures concentrated on fractal sets the exponent can be non-integer. Notice that the formula \eqref{eq:holder-exp} has a similar form to the definition of the local learning coefficient \eqref{eq:llc-formula}. To make this connection precise, we introduce a pseudo-metric induced by the loss function.

\begin{definition}[Loss pseudo-metric]
\label{def:loss-metric}
    The \textbf{loss pseudo-metric} on $\W$ is defined by
    \begin{equation}
        \rho(w,u) = |L(w) - L(u)|.
    \end{equation}
\end{definition}

\begin{remark}
    $\rho$ is only a pseudo-metric as the distance between two distinct points can be zero. 
\end{remark}

Under this pseudo-metric, the sublevel set \eqref{eq:sublevel-set} is precisely the $\epsilon$-ball around $w^*$:
\begin{equation}
    B(w^*, \epsilon) = \{ w \in B(w^*) \mid \rho(w, w^*) < \epsilon \},
\end{equation}
so $V(\epsilon) = \nu(B(w^*, \epsilon))$.

\begin{exercise}[H\"older exponent and learning coefficient] \label{ex:holder-llc}
    Let $\alpha(w^*)$ be the Hölder exponent of the Lebesgue measure on $\W$ at $w^*$ under the loss pseudo-metric. Show that
    \begin{equation*}
        \lambda(w^*) = \alpha(w^*).
    \end{equation*}
\end{exercise}

\subsection{Local learning coefficients of deep linear networks}
\label{sec:llc:dln}

Having established the local learning coefficient as a measure of degeneracy
in the loss landscape, we now compute it explicitly for the two-layer deep
linear network introduced in \cref{sec:prelims}. In \cref{sec:degeneracy},
we saw that DLN parameters exhibit varying degrees of degeneracy depending
on the ranks of the component matrices (\cref{ex:dln-degeneracy},
\cref{rem:rank}). The LLC makes this hierarchy quantitative: different
ranks correspond to different learning coefficients, confirming that
lower-rank parameters are geometrically ``simpler.''

Consider the two-layer DLN $f_{A,B}(x) = BAx$ with
$A \in \R^{H \times M}$ and $B \in \R^{N \times H}$, so the parameter
space is $\W = \R^{(M+N)H}$. Suppose the true input--output relationship
is $y = B_0 A_0 x$ where $A_0 \in \R^{H \times M}$,
$B_0 \in \R^{N \times H}$, and $r = \rank(B_0 A_0)$, and consider the
population loss
\begin{equation}
    L(A, B) = \E_{x \sim q(x)}\!\left[\|BAx - B_0 A_0 x\|^2\right],
\end{equation}
where $q(x)$ has full-rank covariance. The set of minima of this loss function is
$W_0 = \{(A, B) : BA = B_0 A_0\}$.
\citet{Aoyagi+Watanabe2005} derived an exact formula for the learning
coefficient at these minima.

\begin{theorem}[\citealp{Aoyagi+Watanabe2005}]
\label{thm:aoyagi-watanabe}
    For the two-layer DLN $f_{A,B}(x) = BAx$ with
    $A \in \R^{H \times M}$, $B \in \R^{N \times H}$, and true parameter
    $(A_0, B_0)$ with $r = \rank(B_0 A_0)$, at any minimum $w^* \in W_0$ the local learning coefficient $\lambda$ and its multiplicity $\mu$ are given by the following cases:
    \begin{enumerate}
        \item If $N + r \leq M + H$, \; $M + r \leq N + H$, \; and \; $H + r \leq M + N$:
        \begin{enumerate}
            \item If $M + H + N + r$ is even, then $\mu = 1$ and
            \[
                \lambda = \frac{-(H+r)^2 - M^2 - N^2 + 2(H+r)M + 2(H+r)N + 2MN}{8}.
            \]
            \item If $M + H + N + r$ is odd, then $\mu = 2$ and
            \[
                \lambda = \frac{-(H+r)^2 - M^2 - N^2 + 2(H+r)M + 2(H+r)N + 2MN + 1}{8}.
            \]
        \end{enumerate}
        \item If $M + H < N + r$, then $\mu = 1$ and
        $\displaystyle\lambda = \frac{HM - Hr + Nr}{2}$.
        \item If $N + H < M + r$, then $\mu = 1$ and
        $\displaystyle\lambda = \frac{HN - Hr + Mr}{2}$.
        \item If $M + N < H + r$, then $\mu = 1$ and
        $\displaystyle\lambda = \frac{MN}{2}$.
    \end{enumerate}
\end{theorem}

\begin{exercise}[Learning coefficients of two-layer DLNs]
\label{ex:dln-lc}
    Consider the constant-width two-layer DLN from \cref{ex:dln-degeneracy}: $y = BAx$, where $x \in \R^m$ are the inputs, $y \in \R^m$ are the outputs, and $A, B \in \R^{m \times m}$ are linear transformations. The parameter space is $\W = \R^{2m^2}$ and the parameters are the pair $(A, B) \in \W$.
    \begin{enumerate}
        \item Aoyagi \& Watanabe do not assume the network is constant-width, like we do. Simplify their formula for $\lambda$ using this assumption.
        \item How does the learning coefficient $\lambda$ depend on $r$, the rank of $B_0 A_0$? Compare with the results of \cref{ex:dln-degeneracy:general-rank}.
        \item How does $\lambda$ compare with the effective parameter count $d/2$ of a regular model? What does this tell us about the role of singularities in the two-layer DLN?
\end{enumerate}
\end{exercise}


\clearpage
\section{Degeneracy and Bayesian deep learning}
\label{sec:fef}

In \cref{sec:degeneracy,sec:llc}, we studied the geometry of degeneracy in parameter space and introduced the local learning coefficient $\lambda(w^*)$ as a measure of the degree of degeneracy at a local minimum~$w^*$. We now show that this geometry has consequences for learning (in the Bayesian setting). In particular, it creates a trade-off between model fit and model complexity that drives learning, providing a mechanism for \emph{internal model selection}.

\subsection{Watanabe's free energy formula}

Recall the Bayesian learning setup from \cref{sec:bayesian-dl}. We have a parameter--distribution map $\Psi : \W \to \D$ sending each parameter $w$ to a conditional distribution with density $p(y \mid x, w)$, a smooth positive prior $\varphi$ on $\W$, and a data set of $n$ i.i.d.\ examples from a data distribution $q$. Take the loss function to be the negative log-likelihood
\begin{equation*}
    L_n(w) = -\frac{1}{n}\sum_{i=1}^n \log p(y^{(i)} \mid x^{(i)}, w),
    \qquad
    L(w) = \E_{(x,y) \sim q}\!\left[-\log p(y \mid x, w)\right].
\end{equation*}
The posterior $\pi_n$ is given by Bayes' rule \eqref{eq:posterior} and as we saw in \cref{ex:posterior-free-energy,ex:partition-decomposition} it concentrates around regions of parameter space with the lowest local free energy $F_n(\U)$ (\cref{def:local-free-energy}).

We now state one of the main results of SLT: an asymptotic expansion of $F_n(\U)$ as $n \to \infty$. Under certain technical assumptions on the statistical model \citetext{\citealp{Watanabe2018}; see \citealp[][\S 2.3.1]{Lau2025} for a summary}, we have the following theorem.


\begin{theorem}[Local free energy formula]
\label{thm:fef}
    Suppose $\U \subseteq \W$ contains at least one minimiser of $L$, and let $w^* \in \U$ be any such minimiser. Define $\lambda_{\U}$ as the smallest local learning coefficient among all minimisers in $\U$, and let $\mu_{\U}$ be the maximum local multiplicity associated with $\lambda_{\U}$. Then, as $n \to \infty$,
    \begin{equation}
    \label{eq:fef}
        F_n(\U) = nL_n(w^*) + \lambda_{\U} \log n 
              - (\mu_{\U} - 1)\log \log n + O_p(1).
    \end{equation}
\end{theorem}
\begin{remark}
    Since $L_n(w^*)$ is a random variable, so is $F_n(\U)$. The expansion \cref{eq:fef} is an asymptotic statement about random variables where the remainder $O_p(1)$ denotes a term that is bounded in probability as $n \to \infty$.
\end{remark}

To understand what this expansion tells us about Bayesian learning, note that the two leading terms of this expansion have transparent interpretations:
$$F_n \;\approx\;
    \underbrace{n L_n\vphantom{\lambda}}_{\text{data fit}}
    \;+\;
    \underbrace{\lambda \log n}_{\text{complexity}}$$ Let $w_1, w_2 \in \W$ be minimisers of the population loss with $L(w_1) = L(w_2)$ and $\lambda(w_1) \leq \lambda(w_2)$. Then for fixed large $n$, the free energy of a neighbourhood of $w_1$ is smaller than that of $w_2$, so the posterior concentrates around the simpler solution $w_1$.

This is an example of \textbf{internal model selection} where there are two equally good solutions that minimise the population loss. However, our learning algorithm (Bayesian updating) prefers one solution over the other, in particular the one with lower complexity.

\begin{remark}[Regular models and the BIC]
    When the Hessian $\nabla^2 L(w^*)$ is non-degenerate, $\lambda(w^*) = d/2$ and $\mu(w^*) = 1$ by \cref{ex:llc-regular-case}. The free energy
    formula then reduces to $F_n \approx n L_n(w^*) + \frac{d}{2} \log n$,
    recovering the Bayesian Information Criterion from classical statistics.
\end{remark}

\begin{exercise}[Numerical demonstration of the free energy formula]
\label{ex:fef-numerical}
    Consider the statistical model underlying the cubically-parameterised loss from \cref{ex:cubic-parameterisation}. The parameter $\mu \in \R$ indexes the family of normal distributions
        $$p(x \mid \mu) = (2\pi)^{-1/2}\exp\!\left(-\frac{1}{2}(x - \mu^3)^2\right).$$
    Suppose the true distribution is
        $$q(x) = (2\pi)^{-1/2}\exp\!\left(-\frac{1}{2}(x - \mu_0^3)^2\right)$$
    for some fixed $\mu_0$, and equip the model with a standard normal prior
        $$\varphi(\mu) = (2\pi)^{-1/2}\exp\!\left(-\frac12\mu^2\right).$$
    \note{The following exercises are intended to be attempted with a numerical programming framework such as Python or Julia, rather than by hand.}
    %
    \begin{enumerate}
        \item For $\mu_0 = 0$, numerically compute the free energy $F_n$ for many values of $n$ between $n = 100$ and $n = 10{,}000$. Compare the computed values against the asymptotic estimate $F_n \approx nL_n + \lambda \log n$, using the value of $\lambda$ from \cref{ex:cubic-parameterisation:d}.
        \item Repeat part~(a) with $\mu_0 = 5$.
        \item Repeat part~(a) with $\mu_0 = 0.1$. Note that $\mu_0 = 0.1 \neq 0$, so asymptotically $\lambda = 1/2$, as in part~(b). However, observe how for moderate $n$ the free energy behaves more like part~(a) than part~(b). This is another manifestation of the phenomenon from \cref{ex:cubic-parameterisation:e}: the effects of the singularity at $\mu = 0$ extend into a neighbourhood around it.
    \end{enumerate}
\end{exercise}

\subsection{Bayesian phase transitions}

So far we have looked at the consequence of the free energy formula when considering $n$ to be fixed. However, we also find it interesting to see what it tells us about changes in the posterior as $n$ increases. 

For the rest of this section, assume $n$ is sufficiently large that (i) the free energy can be approximated by the first two terms in its expansion, and (ii) $L_n \approx L$ by the law of large numbers, so that the empirical loss can be treated as approximately constant in $n$.

Consider two disjoint regions $\U, \V \subseteq \W$, containing minimisers $u^\ast$ and $v^\ast$ of the population loss $L$ with local learning coefficients $\lambda_{\U}$ and $\lambda_{\V}$ respectively. Applying the free energy formula to each region, the log posterior odds between the two regions satisfy
\begin{equation}
\label{eq:log-odds}
    \log \frac{\pi_n(\U)}{\pi_n(\V)}
    = F_n(\V) - F_n(\U)
    \approx n\,\Delta L_n + \Delta\lambda \cdot \log n
\end{equation}
where $\Delta L_n = L_n(v^\ast) - L_n(u^\ast)$ and $\Delta\lambda = \lambda_{\V} - \lambda_{\U}$.

Consider the case when one region is more accurate but more complex than the other i.e.
\begin{equation*}
    L_n(u^\ast) > L_n(v^\ast)
    \qquad\text{and}\qquad
    \lambda_{\U} < \lambda_{\V}.
\end{equation*}
Then $\Delta L_n < 0$ and $\Delta\lambda > 0$, so the two terms in \eqref{eq:log-odds} have opposite signs. For smaller $n$, the $\Delta\lambda \cdot \log n$ term dominates so the log posterior odds is negative meaning the posterior prefers $\U$ (the simpler, less accurate region). As $n$ increases, the linear term $n\,\Delta L_n$ eventually dominates the logarithmic term, and the posterior shifts to prefer $\V$ (the more complex, more accurate region). Setting the leading terms equal, this change occurs at a critical sample size $n^*$ satisfying
\begin{equation}
\label{eq:critical-n}
    \frac{n^*}{\log n^*} \approx \frac{\Delta\lambda}{\lvert\Delta L_n\rvert}.
\end{equation}
At $n = n^*$, the posterior undergoes a \textbf{Bayesian phase transition}: the region that dominates the posterior changes abruptly. 


\begin{exercise}[Exploring phase transitions]\label{ex:phase-transitions} We show how \emph{phase transitions} can occur when different local free energies change at different rates.
  \begin{enumerate}
    \item Suppose we partition the overall parameter space into a disjoint union $\W = \U \sqcup \V$. Denote the overall free energy by $F_n$ and the local free energies by $F_n(\U)$ and $F_n(\V)$, respectively. Show that the relationship
    \[
    F_n = -\log(e^{-F_n(\U)} + e^{-F_n(\V)})
    \]
    holds.
    
    \item Suppose $F_n(\U) \approx 0.3n + 20 \log n$ and $F_n(\V) \approx 0.5n + 2 \log n$. Plot the overall free energy $F(n)$. Compare against $\min\{F_ n(U), F_n(\V)\}$. What happens around $n = 570$?
    \item In statistical physics, a phase transition is traditionally defined as a discontinuity (or rapid change, for finite-size systems) in the derivatives of the free energy. Plot $\frac{d}{dn}F_n$ and explain why this justifies calling the phenomenon from b) a phase transition.
    \item Change the coefficients of the terms in b). How does this change when a phase transition occurs? Does a phase transition always occur?
    \end{enumerate}
\end{exercise}

\clearpage
\section{Further readings}
\label{sec:readings}

We conclude by providing several pointers to additional literature and resources for those interested in investigating singular learning theory (SLT) in more detail.

\subsection{Other introductions to singular learning theory}

For alternative introductions to SLT, see the following.
\begin{itemize}
    \item
        \citet{Wei+2023} ``Deep Learning is Singular, and That's Good,'' a technical position paper surveying some implications of SLT for deep learning.
    \item
        \citet{Carroll2023} ``Distilling singular learning theory,'' a LessWrong sequence introducing Watanabe's free energy formula and discussing an example of a Bayesian phase transition in a small neural network.
    \item 
        Lecture recordings from the \textit{SLT \& Alignment Summit, 2023}. In particular:
        \begin{itemize}
            \item See the ``SLT Low Road'' lectures \citep{Lau+Chen2023} for an outline of the derivation of Watanabe's free energy formula.
            \item See the ``SLT High Road'' lectures \citep{Murfet+Carroll2023} for a discussion of the free energy formula's implications including Bayesian phase transitions.
        \end{itemize}
    \item
        \citet[][Chapter~2, ``Singular Learning Theory Background'']{Lau2025}, a self-contained technical introduction to the main results of SLT with illustrative examples.
\end{itemize}

See \citet{Furman2024} for a list of mathematical exercises. There is some overlap with exercises included in this tutorial, but there are also several additional exercises.

For a more in-depth introduction to the theoretical foundations of SLT, see Watanabe's two research monographs.
\begin{itemize}
    \item
        \citet{Watanabe2009} ``Algebraic Geometry and Statistical Learning Theory,'' colloquially known as ``the grey book.'' Derives the free energy formula in the realisable case along with other results concerning generalisation properties of Bayesian inference and maximum a posteriori inference.
    \item 
        \citet{Watanabe2018} ``Mathematical Theory of Bayesian Statistics,'' colloquially known as ``the green book.'' An alternative presentation of the free energy formula generalised to the non-realisable case, among other results. Compared to the grey book, the green book has an updated presentation of the main results, but does not contain all of the details of the proofs of the main results from the grey book.
\end{itemize}

While the theoretical foundation of SLT is described across many research papers by Watanabe and others, the main results are collected in self-contained form in these monographs. Watanabe offers a set of lecture slides \citep{Watanabe2023} which may serve as a useful guide to the ``big picture'' while working through the details in the books.
Other key papers include \citet{Watanabe2007,Watanabe2013}. See \citet{Watanabe2024survey} for a more detailed survey.

\clearpage

\subsection{Recent work on singular deep learning}

Over the last few years, a community of researchers have pursued the application of SLT to advancing the science and safety of deep learning.
Some of the ideas behind this research are discussed by \citet{Hoogland+2023,Skalse+2023,PepinLehalleur+2025,Furman2026}.
We briefly survey some key topics in this emerging literature.

\paragraph{Characterising and estimating degeneracy in practice.}
We have seen examples of parameter--function map degeneracy in simple neural networks. Symmetries of neural network parameter--function maps have long been studied, however often the emphasis has been on discrete or globally continuous symmetries rather than additional symmetries localised to subsets of parameter space \citep[see][\S 2.3, for a survey]{FarrugiaRoberts2022}.
For two-layer hyperbolic tangent networks, \citet{FarrugiaRoberts2024} characterises the regions of parameter space which display additional degeneracy, and \citet{FarrugiaRoberts2023} characterises degenerate directions in the parameter--function map.

As discussed, analytically calculating the (local) learning coefficient for deep neural networks is challenging. However, precise formulas have been derived for multi-layer deep linear networks \citep{Aoyagi+Watanabe2005,Aoyagi2024}, two-layer hyperbolic tangent networks \citep{Aoyagi2009}, and certain other architectures. These results assume data is generated from a known ``teacher'' model.

In practice, we lack knowledge of the true data generating process, and we use more complex architectures. Much work has therefore built on the foundational methods of estimating the local learning coefficient with scalable Markov chain Monte Carlo methods \citep{Lau+2025,Hitchcock+Hoogland2025}.
For a practical introduction to learning coefficient estimation, see \citet{Furman2023}.
\citet{Chen+Murfet2025} characterises the sensitivity of local learning coefficient estimation to patterns in sequence models.

\paragraph{Bayesian phase transitions and stagewise development.}
As we have discussed, Watanabe's free energy formula suggests Bayesian deep learning should undergo Bayesian phase transitions under certain conditions. \citet{Carroll2021} studied Bayesian phase transitions in small ReLU networks, and \citet{Chen+2023} studied Bayesian phase transitions in a small feature autoencoder (the ``toy model of superposition'' from \citealp{Elhage+2022}).

In practice, we use stochastic gradient-based optimisation, rather than Bayesian learning, to train neural networks. However, the Bayesian case serves as a model system from which we can derive empirically testable predictions. \citet{Chen+2023} formulate the \emph{Bayesian antecedent hypothesis}, the empirical conjecture that observed phase transitions in trained neural networks correspond to Bayesian phase transitions modelled by Watanabe's free energy formula.
\citet{Chen+2023} study such \emph{dynamical phase transitions} in their toy autoencoder and observe a temporal correspondence between phase changes and estimated LLC increases consistent with the free energy formula.

\citet{Wang+2024,Hoogland+2025} scale this methodology to transformers trained on natural language, finding similar \emph{stagewise development} phenomena with changes in behaviour and internal structure accompanied by LLC changes. \citet{Panickssery+Vaintrob2023,Hoogland+2025,Carroll+2025,Urdshals+Urdshals2025} also study developmental stages in transformers trained on synthetic data.
\citet{Elliott+2026} extends this study to a case of goal misgeneralisation in deep reinforcement learning.

\paragraph{Degeneracy, interpretability, and patterning.}
The LLC provides a single number representing the effective dimensionality of a model. We can derive from the same principles---asymptotic properties of the posterior that reflect degeneracies in the model---more fine-grained tools for probing the internal computational structures of neural networks and how they depend on data.
\begin{itemize}
    \item\textit{Weight- and data-refined LLCs:}
        \citet{Wang+2025differentiation} compute LLCs of individual transformer modules (e.g., different attention heads) or with respect to different subsets of a data set (e.g., natural language versus code). Observing these refined quantities over training reveals modules developing different structures specialising to different kinds of data.
    
    \item\textit{Susceptibilities:}
        Drawing inspiration from electromagnetic susceptibilities in physics
            \citet{Baker+2025,Wang+2025embryology,Gordon+2026}
        develop and apply a methodology for computing loss susceptibilities so as to reveal more fine-grained information about how model internals relate to data.
        
    \item\textit{Bayesian influence functions:}
        Similarly, drawing inspiration from training data attribution in statistics, 
            \citet{Kreer+2025,Lee+2025,Adam+2025}
        develop a Bayesian generalisation of classical influence functions that allows the influence functions to be sensitive to higher-order degeneracy in the model.
\end{itemize}
Weight- and data-refined LLCs are themselves LLCs with a different model or data set, and so the same scalable Markov chain Monte Carlo methods can be used to estimate them as for LLCs. Moreover, like the LLC, susceptibilities and Bayesian influence functions can be approximated as expectations over a localised posterior distribution, and so similar estimation methods can be used for these quantities too.

\citet{Wang+Murfet2026} develop a methodology, \emph{patterning}, for making targeted changes to the training distribution so as to elicit certain changes in the development of neural network internal structure or generalisation behaviour.

\paragraph{Foundations of singular deep learning.}

There has been some work on developing the foundational theory of SLT and deep learning.
%
For example, \citet{Elliott+2026} generalise Watanabe's free energy formula from Bayesian inference to a generalised non-stationary energy-based inference setting, so as to account for stagewise development in deep reinforcement learning.
%
Beyond the setting of Bayesian inference, the general role degeneracy plays in the learning dynamics of stochastic gradient-based optimisation remains to be characterised.

Finally, there has been some attempt to theoretically investigate the links between degeneracy in deep learning and computational structure in models.
%
\citet{Clift+2021,Waring2021,Xu2021,Murfet2024,Murfet+Troiani2025} study degeneracies in a statistical model based on a parameterisation of the space of Turing machine programs, exploring links between degeneracy and computational structure.
% 
\citet[][Chapter~5]{Lau2025} and \citet{Urdshals+2025} develop a theory of minimum description length in degenerate statistical models.


% % % 
% References

\clearpage
\bibliography{biblo}
\addcontentsline{toc}{section}{References}


% % %
% List of solutions
\begin{solutionsonly}
\clearpage
\appendix
\section{Solutions}
\label{apx:solutions}


% command for spacing between solution sections
\newcommand\solutionlistskip{\protect\addvspace{2ex}}
\end{solutionsonly}

% \clearpage
% \subsection[Solutions to exercises from Section~1]{Solutions to exercises from \cref{sec:prelims}}

\begin{solution}[ex:biased-neurons]
For each architecture, we add a bias vector to each layer's output before the next transformation (or activation function) is applied.

\begin{enumerate}
    \item \textbf{A linear neuron with bias} (\cref{eg:neuron}).
    Let $\W = \R^{d+1}$, encoding a weight vector $w \in \R^d$ and a scalar bias $b \in \R$. The parameter--function map sends $(w, b)$ to $f_{w,b} : \R^d \to \R$ defined by
    \begin{equation*}
        f_{w,b}(x) = w^\top x + b.
    \end{equation*}

    \item \textbf{Multi-linear neural network with bias} (\cref{eg:linear}).
    Let $\W = \R^{m^2 + m}$, encoding a matrix $W \in \R^{m \times m}$ and a bias vector $b \in \R^m$. The parameter--function map sends $(W, b)$ to $f_{W,b} : \R^m \to \R^m$ defined by
    \begin{equation*}
        f_{W,b}(x) = Wx + b.
    \end{equation*}

    \item \textbf{Deep linear network with bias} (\cref{eg:dln}).
    Let $\W = \R^{2mh + h + m}$, encoding matrices $A \in \R^{h \times m}$, $B \in \R^{m \times h}$ and bias vectors $b_A \in \R^h$, $b_B \in \R^m$. The parameter--function map sends $(A, b_A, B, b_B)$ to $f_w : \R^m \to \R^m$ defined by
    \begin{equation*}
        f_w(x) = B(Ax + b_A) + b_B.
    \end{equation*}

    \item \textbf{Multi-layer perceptron with bias} (\cref{eg:mlp}).
    The parameter space is $\W = \R^{2mh + h + m}$ as above. The parameter--function map sends $(A, b_A, B, b_B)$ to $f_w : \R^m \to \R^m$ defined by
    \begin{equation*}
        f_w(x) = B \cdot \sigma(Ax + b_A) + b_B. \qedhere
    \end{equation*}
\end{enumerate}
\end{solution}

\begin{solution}[ex:empirical-population-loss]
\begin{enumerate}
    \item
        Since the data pairs $(x^{(i)}, y^{(i)})$ are sampled independently and
        identically from $q$, each per-example loss $L_{x^{(i)},y^{(i)}}(w)$ is
        an identically distributed random variable with expectation
        \begin{equation*}
            \E_{(x,y)\sim q}[L_{x,y}(w)] = L(w).
        \end{equation*}
        Therefore, by linearity of expectation,
        \begin{equation*}
            \E[L_n(w)]
            = \E\left[\frac{1}{n}\sum_{i=1}^n L_{x^{(i)},y^{(i)}}(w)\right]
            = \frac{1}{n}\sum_{i=1}^n \E[L_{x^{(i)},y^{(i)}}(w)]
            = \frac{1}{n} \cdot n \cdot L(w)
            = L(w).
        \end{equation*}

    \item
        The random variables $L_{x^{(1)},y^{(1)}}(w), L_{x^{(2)},y^{(2)}}(w),
        \ldots$ are independent and identically distributed with common mean
        $L(w)$. By the strong law of large numbers,
        \begin{equation*}
            L_n(w)
            = \frac{1}{n}\sum_{i=1}^n L_{x^{(i)},y^{(i)}}(w)
            \xrightarrow{n\to\infty} L(w)
        \end{equation*}
        almost surely. \qedhere
\end{enumerate}
\end{solution}

\begin{solution}[ex:nll-mse]
\begin{enumerate}
    \item
        Since labels are sampled independently, the likelihood factorises as a
        product of Gaussian densities:
        \begin{align*}
            p(Y_n \mid X_n, w)
            &= \prod_{i=1}^n p(y^{(i)} | x^{(i)}, w)
            \\
            &= \prod_{i=1}^n
                \frac{1}{(2\pi)^{m/2}}
                \exp\!\left(
                    -\frac{\|y^{(i)} - f_w(x^{(i)})\|^2}{2}
                \right).
        \end{align*}

    \item
        Taking the negative logarithm of the above expression,
        \begin{align*}
            -\log p(Y_n \mid X_n, w)
            &= -\sum_{i=1}^n \left[
                -\frac{m}{2}\log(2\pi)
                - \frac{\|y^{(i)} - f_w(x^{(i)})\|^2}{2}
            \right]
            \\
            &= \frac{nm}{2}\log(2\pi)
                + \frac{1}{2}\sum_{i=1}^n
                    \|y^{(i)} - f_w(x^{(i)})\|^2.
        \end{align*}

    \item
        From the mean squared error definition \eqref{eq:mse},
        \begin{equation*}
            -\log p(Y_n \mid X_n, w)
            = \frac{nm}{2}\log(2\pi) + \frac{n}{2} L_n(w).
        \end{equation*}
        The right-hand side is a strictly increasing affine function of
        $L_n(w)$ (with positive coefficient $n/2$), and the additive constant
        $\frac{nm}{2}\log(2\pi)$ is independent of $w$. Therefore,
        \begin{equation*}
            \argmax_{w\in\W} p(Y_n \mid X_n, w)
            = \argmin_{w\in\W} \left(-\log p(Y_n \mid X_n, w)\right)
            = \argmin_{w\in\W} L_n(w). \qedhere
        \end{equation*}
\end{enumerate}
\end{solution}

\begin{solution}[ex:nll-cross-entropy]
\begin{enumerate}
    \item
        Since labels are sampled independently, the likelihood factorises as
        \begin{equation*}
            p(C_n \mid X_n, w)
            = \prod_{i=1}^n p(c^{(i)} | x^{(i)}, w)
            = \prod_{i=1}^n f_w(c^{(i)} \mid x^{(i)}).
        \end{equation*}

    \item
        Taking the negative logarithm,
        \begin{equation*}
            -\log p(C_n \mid X_n, w)
            = -\sum_{i=1}^n \log f_w(c^{(i)} \mid x^{(i)}).
        \end{equation*}

    \item
        Since $y^{(i)} = \delta_{c^{(i)}}$ is the Dirac distribution on class
        $c^{(i)}$, we have $y^{(i)}(j) = 1$ if $j = c^{(i)}$ and $y^{(i)}(j)
        = 0$ otherwise. Substituting into the cross entropy loss
        \eqref{eq:cross-entropy},
        \begin{equation*}
            L_{x^{(i)},y^{(i)}}(w)
            = -\sum_{j=1}^m y^{(i)}(j) \log f_w(j \mid x^{(i)})
            = -\log f_w(c^{(i)} \mid x^{(i)}).
        \end{equation*}
        Therefore the empirical loss is
        \begin{equation*}
            L_n(w)
            = \frac{1}{n}\sum_{i=1}^n L_{x^{(i)},y^{(i)}}(w)
            = -\frac{1}{n}\sum_{i=1}^n \log f_w(c^{(i)} \mid x^{(i)})
            = \frac{1}{n}\left(-\log p(C_n \mid X_n, w)\right).
        \end{equation*}
        Since the negative log-likelihood is $n \cdot L_n(w)$, a positive
        multiple of the empirical loss, we conclude
        \begin{equation*}
            \argmax_{w\in\W} p(C_n \mid X_n, w)
            = \argmin_{w\in\W} L_n(w). \qedhere
        \end{equation*}
\end{enumerate}
\end{solution}

\begin{solution}[ex:nll-energy]
\begin{enumerate}
    \item
        Since labels are sampled independently, the likelihood factorises as
        \begin{equation*}
            p(Y_n \mid X_n, w)
            = \prod_{i=1}^n p(y^{(i)} | x^{(i)}, w)
            = \prod_{i=1}^n
                \frac
                    {\exp\!\left(-L_{x^{(i)},y^{(i)}}(w)\right)}
                    {Z(x^{(i)}, w)}.
        \end{equation*}
        Taking the negative logarithm,
        \begin{align*}
            -\log p(Y_n \mid X_n, w)
            &= -\sum_{i=1}^n \left[
                -L_{x^{(i)},y^{(i)}}(w) - \log Z(x^{(i)}, w)
            \right]
            \\
            &= \sum_{i=1}^n L_{x^{(i)},y^{(i)}}(w)
                + \sum_{i=1}^n \log Z(x^{(i)}, w).
        \end{align*}

    \item
        Suppose $L_{x,y}(w) = \ell(f_w(x) - y)$. Making the substitution
        $u = f_w(x) - z$ (so that $dz = du$) in the definition of the
        partition function,
        \begin{equation*}
            Z(x, w)
            = \int_{\R^m} \exp\!\left(-\ell(f_w(x) - z)\right) dz
            = \int_{\R^m} \exp\!\left(-\ell(u)\right) du.
        \end{equation*}
        The right-hand side is a constant independent of both $x$ and $w$.
        Therefore, by part~(a), the negative log-likelihood is
        \begin{equation*}
            -\log p(Y_n \mid X_n, w)
            = \sum_{i=1}^n L_{x^{(i)},y^{(i)}}(w) + C
            = n \cdot L_n(w) + C
        \end{equation*}
        where $C = n \log \int_{\R^m} \exp(-\ell(u))\,du$ is independent of
        $w$. Since the negative log-likelihood is an affine function of
        $L_n(w)$ with positive coefficient,
        \begin{equation*}
            \argmax_{w\in\W} p(Y_n \mid X_n, w)
            = \argmin_{w\in\W} L_n(w). \qedhere
        \end{equation*}
\end{enumerate}
\end{solution}

\begin{solution}[ex:gibbs]
    By definition of the negative log-likelihood,
    \begin{equation*}
        n L_n(w)
        = -\sum_{i=1}^{n} \log p(y^{(i)} \mid x^{(i)}, w).
    \end{equation*}
    Exponentiating both sides,
    \begin{equation*}
        \exp\!\big({-}n L_n(w)\big)
        = \exp\!\left(\sum_{i=1}^{n} \log p(y^{(i)} \mid x^{(i)}, w)\right)
        = \prod_{i=1}^{n} p(y^{(i)} \mid x^{(i)}, w).
    \end{equation*}
    Substituting into the definitions of the posterior
    \eqref{eq:posterior} and partition function \eqref{eq:partition-function}
    gives the stated expressions.
\end{solution}

\begin{solution}[ex:posterior-free-energy]
\begin{enumerate}
    \item
        Integrating the posterior \eqref{eq:posterior} over $\U$,
        \begin{equation*}
            \pi_n(\U)
            = \int_{\U} \pi_n(w) \, dw
            = \int_{\U} \frac{1}{Z_n} \varphi(w)
                \prod_{i=1}^{n} p(y^{(i)} \mid x^{(i)}, w) \, dw
            = \frac{Z_n(\U)}{Z_n}.
        \end{equation*}

    \item
        By part~(a),
        \begin{equation*}
            \log \frac{\pi_n(\U)}{\pi_n(\V)}
            = \log \frac{Z_n(\U) / Z_n}{Z_n(\V) / Z_n}
            = \log Z_n(\U) - \log Z_n(\V)
            = F_n(\V) - F_n(\U). \qedhere
        \end{equation*}
\end{enumerate}
\end{solution}

\begin{solution}[ex:partition-decomposition]
\begin{enumerate}
    \item
        Since $\U_1, \ldots, \U_k$ are disjoint and cover $\W$, the integral
        over $\W$ decomposes as a sum of integrals over each $\U_j$:
        \begin{align*}
            Z_n
            &= \int_{\W} \varphi(w) \exp\!\big({-}n L_n(w)\big)\, dw
            \\
            &= \sum_{j=1}^{k}
                \int_{\U_j} \varphi(w) \exp\!\big({-}n L_n(w)\big)\, dw
            = \sum_{j=1}^{k} Z_n(\U_j).
        \end{align*}
        By definition, $Z_n(\U_j) = e^{-F_n(\U_j)}$, so
        $Z_n = \sum_{j=1}^{k} e^{-F_n(\U_j)}$.
        Taking $-\log$ of both sides,
        \begin{equation*}
            F_n = -\log Z_n = -\log \sum_{j=1}^{k} e^{-F_n(\U_j)}.
        \end{equation*}

    \item
        Since $F_n(\U_1) < F_n(\U_j)$ for all $j \neq 1$, the corresponding
        partition functions satisfy $Z_n(\U_1) > Z_n(\U_j)$ for all $j \neq 1$.
        The lower bound on $Z_n$ is immediate:
        \begin{equation*}
            Z_n = \sum_{j=1}^{k} Z_n(\U_j) > Z_n(\U_1).
        \end{equation*}
        For the upper bound, since each $Z_n(\U_j) < Z_n(\U_1)$,
        \begin{equation*}
            Z_n = \sum_{j=1}^{k} Z_n(\U_j) < k \cdot Z_n(\U_1).
        \end{equation*}
        Taking $-\log$ (which reverses the inequalities) yields
        \begin{equation*}
            F_n(\U_1) - \log k < F_n < F_n(\U_1).
        \end{equation*}
        The overall free energy is therefore determined by the region with
        the lowest local free energy, up to a correction of at most $\log k$.
        In particular, if the gap $F_n(\U_j) - F_n(\U_1)$ grows with $n$ while
        $k$ remains fixed, then $F_n \to F_n(\U_1)$.
        \qedhere
\end{enumerate}
\end{solution}

% \clearpage
% \subsection[Solutions to exercises from Section~2]{Solutions to exercises from \cref{sec:degeneracy}}
% \addtocontents{loe}{\solutionlistskip}

\begin{solution}[ex:univariate-degeneracy]
\begin{enumerate}
    \item Since $\W = \R$ is one-dimensional, the directional derivative in direction $v = 1$ reduces to the ordinary derivative:
    \begin{equation*}
        \frac{\partial \Phi}{\partial w} = \frac{d}{dw} w = 1.
    \end{equation*}

    \item The directional derivative is $1 \neq 0$ for all $w \in \W$. Therefore, the parameter--function map is \emph{not degenerate} at any point. \qedhere
\end{enumerate}
\end{solution}

\begin{solution}[ex:cubic-degeneracy]
\begin{enumerate}
    \item The hypothesis class of \cref{ex:univariate-degeneracy} is $\F = \{f_w = w : w \in \R\} = \R$. The hypothesis class here is $\F = \{f_w = w^3 : w \in \R\}$. Since $w \mapsto w^3$ is a bijection on $\R$, as $w$ ranges over $\R$, $w^3$ takes every real value exactly once. So $\F = \R$ in both cases.

    \item The directional derivative in direction $v = 1$ is the ordinary derivative:
    \begin{equation*}
        \frac{\partial \Phi}{\partial w} = \frac{d}{dw} w^3 = 3w^2.
    \end{equation*}

    \item The directional derivative $3w^2 = 0$ if and only if $w = 0$. So the parameter--function map is degenerate at $w = 0$ only. It is somewhere degenerate but not everywhere degenerate. \qedhere
\end{enumerate}
\end{solution}

\begin{solution}[ex:product-degeneracy]
\begin{enumerate}
    \item The hypothesis class is $\F = \{f_w = ab : (a, b) \in \R^2\} = \R$, since for any target $c \in \R$ we can choose, e.g., $a = c$ and $b = 1$. This is the same as in \cref{ex:univariate-degeneracy}.

    \item The directional derivative in direction $v = (1, 0)$ is the partial derivative with respect to $a$:
    \begin{equation*}
        \frac{\partial \Phi}{\partial a} = \frac{\partial}{\partial a}(ab) = b.
    \end{equation*}

    \item The directional derivative in direction $v = (0, 1)$ is the partial derivative with respect to $b$:
    \begin{equation*}
        \frac{\partial \Phi}{\partial b} = \frac{\partial}{\partial b}(ab) = a.
    \end{equation*}

    \item From parts (b) and (c), the direction $(1, 0)$ is degenerate at $(a, b)$ if and only if $b = 0$, and the direction $(0, 1)$ is degenerate if and only if $a = 0$. So the parameter--function map is degenerate in at least one of these coordinate directions precisely on the set $\{(a, b) : a = 0 \text{ or } b = 0\}$, i.e.\ the union of the two coordinate axes.

    \item The directional derivative in a general direction $v = (v_1, v_2)$ is proportional to
    \begin{equation*}
        v_1 \frac{\partial \Phi}{\partial a} + v_2 \frac{\partial \Phi}{\partial b} = v_1 b + v_2 a.
    \end{equation*}
    This is a single linear equation in two unknowns $(v_1, v_2)$.
    \begin{itemize}
        \item If $a = b = 0$: the equation $0 = 0$ is trivially satisfied, so every direction is degenerate.
        \item If $(a, b) \neq (0, 0)$: the coefficient vector $(b, a)$ is nonzero, so the solution space is one-dimensional. A nonzero solution is $v = (-a, b)$: indeed, $(-a) b + b \cdot a = 0$.
    \end{itemize}
    Therefore, the parameter--function map is degenerate \emph{everywhere}.
    Note that the previous part found degeneracy only on the coordinate axes
    because it checked only the coordinate directions. Checking for
    parameter--function map degeneracy requires checking all possible
    directions.
    \qedhere
\end{enumerate}
\end{solution}

\begin{solution}[ex:sum-degeneracy]
\begin{enumerate}
    \item The hypothesis class is $\F = \{f_w = a + b : (a, b) \in \R^2\} = \R$, since for any target $c \in \R$ we can choose, e.g., $a = c$ and $b = 0$. This is the same as in \cref{ex:univariate-degeneracy}.

    \item The directional derivative in direction $v = (1, -1)$ is
    \begin{equation*}
        1 \cdot \frac{\partial}{\partial a}(a + b) + (-1) \cdot \frac{\partial}{\partial b}(a + b) = 1 - 1 = 0.
    \end{equation*}

    \item The direction $v = (1, -1)$ gives a zero directional derivative at \emph{every} point $(a, b) \in \W$. Therefore, the parameter--function map is everywhere degenerate. This reflects the fact that the map $(a, b) \mapsto a + b$ has a global continuous symmetry: translating along the direction $(1, -1)$ does not change the output. \qedhere
\end{enumerate}
\end{solution}

\begin{solution}[ex:sum-symmetry]
\begin{enumerate}
    \item The transformation $T_t$ translates the parameter $(a, b)$ by $t$ in the direction $(1, -1)$. The orbit of any parameter $(a, b)$ under this family is the line $\{(a + t, b - t) : t \in \R\}$, which has slope $-1$ in the $(a, b)$-plane.

    \item We verify the three conditions of a continuous symmetry.
    \begin{enumerate}[label=(\roman*)]
        \item Symmetry: $\Phi(T_t(a, b)) = (a + t) + (b - t) = a + b = \Phi(a, b)$ for all $(a, b)$ and $t$.
        \item Identity: $T_0(a, b) = (a + 0, b - 0) = (a, b)$.
        \item Differentiability: $t \mapsto T_t(a, b) = (a + t, b - t)$ is linear in $t$, hence differentiable.
    \end{enumerate}

    \item We have $\frac{d}{dt} T_t(a, b) = (1, -1)$ for all $t$. In particular, $\left.\frac{d}{dt}\right|_{t=0} T_t(a, b) = (1, -1) \neq (0, 0)$ for every $(a, b) \in \W$. So the symmetry is non-trivial at every parameter. \qedhere
\end{enumerate}
\end{solution}

\begin{solution}[ex:relu-scaling]
\begin{enumerate}
    \item We consider three cases. If $z > 0$, then $\alpha z > 0$ (since $\alpha > 0$), so $\relu(\alpha z) = \alpha z = \alpha \relu(z)$. If $z = 0$, then $\relu(\alpha \cdot 0) = 0 = \alpha \cdot 0 = \alpha \relu(z)$. If $z < 0$, then $\alpha z < 0$, so $\relu(\alpha z) = 0 = \alpha \cdot 0 = \alpha \relu(z)$.

    \item We verify the three conditions.
    \begin{enumerate}[label=(\roman*)]
        \item Symmetry: using positive homogeneity with $\alpha = e^t > 0$,
        \begin{equation*}
            f_{T_t(a,b)}(x)
            = e^{-t}b \cdot \relu(e^t a x)
            = e^{-t}b \cdot e^t \relu(ax)
            = b \cdot \relu(ax)
            = f_{a,b}(x).
        \end{equation*}
        \item Identity: $T_0(a, b) = (e^0 a,\, e^0 b) = (a, b)$.
        \item Differentiability: $t \mapsto (e^t a,\, e^{-t} b)$ is smooth.
    \end{enumerate}

    \item The degenerate direction at $(a, b)$ is
    \begin{equation*}
        \left.\frac{d}{dt}\right|_{t=0} T_t(a, b)
        = \left.\frac{d}{dt}\right|_{t=0} (e^t a,\, e^{-t} b)
        = (a, -b).
    \end{equation*}
    Since $(a, b) \neq (0, 0)$, we have $(a, -b) \neq (0, 0)$, so this is a non-zero degenerate direction.

    \item At $(a, b) = (0, 0)$, we have $\left.\frac{d}{dt}\right|_{t=0} T_t(0, 0) = (0, 0)$, so the symmetry is trivial.

    The parameter--function map is nonetheless degenerate at the origin. Both partial derivatives of $\Phi$ vanish there:
    \begin{align*}
        \frac{\partial}{\partial a} f_{a,b}(x) \bigg|_{(0,0)}
        &= \lim_{\epsilon \to 0}
            \frac{f_{\epsilon, 0}(x) - f_{0,0}(x)}{\epsilon}
        = \lim_{\epsilon \to 0}
            \frac{0 \cdot \relu(\epsilon x)}{\epsilon}
        = 0,
    \\
        \frac{\partial}{\partial b} f_{a,b}(x) \bigg|_{(0,0)}
        &= \relu(0 \cdot x) = 0.
    \end{align*}
    Since $\nabla \Phi(0,0) = 0$, every direction is degenerate at
    the origin. This degeneracy is not explained by the scaling symmetry.
    \qedhere
\end{enumerate}
\end{solution}

\begin{solution}[ex:rotation-symmetry]
\begin{enumerate}
    \item For orthogonal $R$ (i.e.\ $R^\top R = I$),
    \begin{equation*}
        f_{RW}(x) = (RW)^\top(RW)x = W^\top R^\top R\, W x = W^\top W x = f_W(x).
    \end{equation*}

    \item We verify the three conditions.
    \begin{enumerate}[label=(\roman*)]
        \item Symmetry: $R(\theta)$ is orthogonal for all $\theta$, since
        \begin{equation*}
            R(\theta)^\top R(\theta)
            = \begin{pmatrix}
                \cos\theta & \sin\theta \\
                -\sin\theta & \cos\theta
            \end{pmatrix}
            \begin{pmatrix}
                \cos\theta & -\sin\theta \\
                \sin\theta & \cos\theta
            \end{pmatrix}
            = \begin{pmatrix}
                1 & 0 \\ 0 & 1
            \end{pmatrix}.
        \end{equation*}
        So $f_{R(\theta)W} = f_W$ by part~(a).
        \item Identity: $R(0) = I$, so $T_0(W) = W$.
        \item Differentiability: the entries of $R(\theta)W$ are smooth functions of $\theta$ (they involve $\cos\theta$ and $\sin\theta$ multiplied by entries of $W$).
    \end{enumerate}

    \item The degenerate direction at $W$ is
    \begin{equation*}
        \left.\frac{d}{d\theta}\right|_{\theta=0} R(\theta) W = R'(0)\, W.
    \end{equation*}
    We compute
    \begin{equation*}
        R'(\theta)
        = \begin{pmatrix}
            -\sin\theta & -\cos\theta \\
            \cos\theta & -\sin\theta
        \end{pmatrix},
        \qquad
        R'(0)
        = \begin{pmatrix}
            0 & -1 \\
            1 & 0
        \end{pmatrix}.
    \end{equation*}
    Writing $W = \begin{pmatrix} w_1^\top \\ w_2^\top \end{pmatrix}$ where $w_1, w_2 \in \R^m$ are the two rows, the degenerate direction is the matrix
    \begin{equation*}
        R'(0)\, W
        = \begin{pmatrix} -w_2^\top \\ w_1^\top \end{pmatrix},
    \end{equation*}
    which swaps the two rows and negates one. This direction is non-zero whenever $W \neq 0$.

    \item The orthogonal group $O(h)$ consists of all $h \times h$ matrices satisfying $R^\top R = I$. This constraint comprises $\frac{h(h+1)}{2}$ independent scalar equations (the entries of the symmetric matrix $R^\top R$ on and above the diagonal). Since $R$ has $h^2$ entries, the dimension of $O(h)$ is
    \begin{equation*}
        h^2 - \frac{h(h+1)}{2} = \frac{h(h-1)}{2}.
    \end{equation*}
    Each independent direction in $O(h)$ at the identity gives rise to a one-parameter continuous symmetry and hence a degenerate direction at each $W$. These directions are the $h \times h$ skew-symmetric matrices $S$ (satisfying $S^\top = -S$), since differentiating $R(t)^\top R(t) = I$ at $t = 0$ gives $R'(0)^\top + R'(0) = 0$. The space of such matrices has dimension $\frac{h(h-1)}{2}$ (the entries strictly above the diagonal are free, and the rest are determined).

    So $O(h)$ contributes $\frac{h(h-1)}{2}$ independent continuous symmetries. For $h = 2$ this gives $1$, matching the single rotation symmetry from part~(b). \qedhere
\end{enumerate}
\end{solution}

\begin{solution}[ex:localised-symmetry]
\begin{enumerate}
    \item $T^{(a)}_t$ translates along the $a$-axis: it shifts the first coordinate by $t$ while leaving the second fixed. $T^{(b)}_t$ translates along the $b$-axis.

    \item $T^{(a)}_t$ is a symmetry at $(a, b)$ for all $t$ if and only if $f_{a+t,b} = f_{a,b}$ for all $t$, i.e., $(a+t)b = ab$ for all $t$. This simplifies to $tb = 0$ for all $t$, which requires $b = 0$. So $\{T^{(a)}_t\}$ is a symmetry when restricted to the $a$-axis $\{(a, 0) : a \in \R\}$.

    Similarly, $\{T^{(b)}_t\}$ is a symmetry when restricted to the $b$-axis $\{(0, b) : b \in \R\}$.

    \item From \cref{ex:product-degeneracy:axes-directions}, the coordinate directions $(1, 0)$ and $(0, 1)$ are degenerate precisely on the coordinate axes: $(1, 0)$ is degenerate on $\{b = 0\}$ and $(0, 1)$ is degenerate on $\{a = 0\}$. This matches exactly: the axis $\{b = 0\}$ is where the $a$-translation symmetry acts, providing the degenerate direction $(1, 0)$, and the axis $\{a = 0\}$ is where the $b$-translation symmetry acts, providing $(0, 1)$. \qedhere
\end{enumerate}
\end{solution}

\begin{solution}[ex:dln-degeneracy]
\begin{enumerate}
    \item
        Encoding $(A, B)$ and $(\delta\!A, \delta\!B)$ as vectors in the underlying parameter space $\R^{2m^2}$, the directional derivative is
        \begin{equation*}
            D_{(\delta\!A, \delta\!B)} \Phi(A, B)(x)
            = \lim_{\epsilon \to 0}
                \frac{
                    f_{A + \epsilon\delta\!A,\, B + \epsilon\delta\!B}(x)
                    - f_{A,B}(x)
                }{\epsilon}.
        \end{equation*}
        Expanding the matrix product:
        \begin{align*}
            f_{A + \epsilon\delta\!A,\, B + \epsilon\delta\!B}(x)
            &= (B + \epsilon\delta\!B)(A + \epsilon\delta\!A)\,x
        \\  &= BAx
                + \epsilon(B\,\delta\!A + \delta\!B\, A)\,x
                + \epsilon^2 \delta\!B\,\delta\!A\, x.
        \end{align*}
        Subtracting $f_{A,B}(x) = BAx$, dividing by $\epsilon$, and taking $\epsilon \to 0$ gives
        \begin{equation*}
            D_{(\delta\!A, \delta\!B)} \Phi(A, B)(x)
            = (B\,\delta\!A + \delta\!B\, A)\,x. \qedhere
        \end{equation*}

    \item
        At $(A, B) = (0, 0)$, part~(a) gives
            $D_{(\delta\!A, \delta\!B)} \Phi(0, 0)(x)
            = (0 \cdot \delta\!A + \delta\!B \cdot 0)\,x = 0$
        for every $(\delta\!A, \delta\!B)$ and every $x$.
        Every non-zero direction is therefore degenerate.
        The subspace of degenerate directions is the entire parameter space
        $\R^{2m^2}$, which has dimension $2m^2$.

    \item
        By part~(a), $(\delta\!A, \delta\!B)$ is a degenerate direction if and only if $B\,\delta\!A + \delta\!B\, A = 0$, that is, $\delta\!B\, A = -B\,\delta\!A$. Since $A$ is invertible, right-multiplying by $A^{-1}$ gives $\delta\!B = -B\,\delta\!A\, A^{-1}$.
        Since $\delta\!A \in \R^{m \times m}$ is free and $\delta\!B$ is
        uniquely determined, the subspace of degenerate directions has
        dimension $m^2$.

    \item
    \begin{enumerate}
        \item
            The $i$-th row of $\delta\!B\, A$ is
                $(\delta\!B)_i A$,
            a linear combination of the rows of $A$.
            Therefore the image of $T_A$ consists of all matrices whose rows
            lie in the row space of $A$.
            Each of the $m$ rows can be any vector in the
            $r_A$-dimensional row space, so
                $\rank(T_A) = m\,r_A$.

        \item
            Similarly, the $j$-th column of $B\,\delta\!A$ is
                $B\,(\delta\!A)^j$,
            a linear combination of the columns of $B$.
            The image of $T_B$ consists of all matrices whose columns lie in
            the column space of $B$, and
                $\rank(T_B) = m\,r_B$.

        \item
            A matrix lies in
                $\im(T_A) \cap \im(T_B)$
            if and only if its rows lie in the row space of $A$ and its
            columns lie in the column space of $B$.
            Such a matrix can be written as $U \Lambda V$ where
            $U \in \R^{m \times r_B}$ has columns spanning the column space
            of $B$ and $V \in \R^{r_A \times m}$ has rows spanning the row
            space of $A$, with $\Lambda \in \R^{r_B \times r_A}$ free.
            The dimensionality of this intersection is therefore $r_A\,r_B$.

        \item
            The image of $T_D$ is
                $\im(T_D) = \im(T_B) + \im(T_A)$,
            since for any $(\delta\!A, \delta\!B)$ we have
                $T_D(\delta\!A, \delta\!B) = T_B(\delta\!A) + T_A(\delta\!B)$,
            and conversely any element of $\im(T_B) + \im(T_A)$ can be
            realised by choosing appropriate $\delta\!A$ and $\delta\!B$.
            By Grassmann's identity,
            \begin{align*}
                \rank(T_D)
                &= \dim(\im(T_A) + \im(T_B))
            \\  &= \rank(T_A) + \rank(T_B)
                    - \dim(\im(T_A) \cap \im(T_B))
            \\  &= m\,r_A + m\,r_B - r_A\,r_B.
            \end{align*}

        \item
            By the rank--nullity theorem, the nullity of $T_D$ is
            \begin{align*}
                &\mathrel{\phantom{=}} 2m^2 - (m\,r_A + m\,r_B - r_A\,r_B)
            \\  &= 2m^2 - m\,r_A - m\,r_B + r_A\,r_B
            \\  &= m^2 + (m - r_A)(m - r_B).
            \end{align*}
            The kernel of $T_D$ is exactly the set of $(\delta\!A, \delta\!B)$
            such that $B\,\delta\!A + \delta\!B\, A = 0$, which by part~(a)
            is precisely the subspace of degenerate directions. \qedhere
    \end{enumerate}
\end{enumerate}
\end{solution}

\begin{solution}[ex:mlp-degeneracy]
\begin{enumerate}
    \item
        At $a_i = 0$, the parameter--function map reduces to
        \begin{equation*}
            f_w(x)
            = d + 0 \cdot \sigma(b_i x + c_i) + \sum_{j \neq i} a_j\,\sigma(b_j x + c_j)
            = d + \sum_{j \neq i} a_j\,\sigma(b_j x + c_j).
        \end{equation*}
        This expression does not involve $b_i$ or $c_i$ at all. Therefore $\Phi$ is constant in the $b_i$ and $c_i$ directions at this point, so $D_{\delta b_i}\Phi(w) = 0$ and $D_{\delta c_i}\Phi(w) = 0$. Both are degenerate directions.

        Note that this argument is purely algebraic, following from the multiplicative structure $a_i \cdot \sigma(\ldots)$, and requires no conditions on~$\sigma$ (not even continuity or differentiability).

    \item
        At $b_i = 0$, unit~$i$ computes the constant $\sigma(0 \cdot x + c_i) = \sigma(c_i)$ for all $x$, so
        \begin{equation*}
            f_w(x) = d + a_i\,\sigma(c_i) + \sum_{j \neq i} a_j\,\sigma(b_j x + c_j).
        \end{equation*}
        In the direction $(\delta a_i, \delta d) = (1, -\sigma(c_i))$, we compute
        \begin{align*}
            f_{w + \epsilon\,v}(x)
            &= (d - \epsilon\,\sigma(c_i))
                + (a_i + \epsilon)\,\sigma(c_i)
                + \sum_{j \neq i} a_j\,\sigma(b_j x + c_j)
        \\  &= d + a_i\,\sigma(c_i)
                + \epsilon\bigl(\sigma(c_i) - \sigma(c_i)\bigr)
                + \sum_{j \neq i} a_j\,\sigma(b_j x + c_j)
            = f_w(x).
        \end{align*}
        So $D_{(\delta a_i, \delta d)}\Phi(w) = 0$ and this direction is degenerate. Intuitively, increasing $a_i$ scales up the constant contribution $a_i\,\sigma(c_i)$ to the output, and decreasing $d$ by the same amount compensates exactly.

    \item When $(b_i, c_i) = (b_j, c_j)$, units $i$ and $j$ compute the same activation, so their combined contribution to $f_w$ is
        \begin{equation*}
            a_i\,\sigma(b_i x + c_i) + a_j\,\sigma(b_j x + c_j)
            = (a_i + a_j)\,\sigma(b_i x + c_i).
        \end{equation*}
        In the direction $(\delta a_i, \delta a_j) = (1, -1)$, this combined contribution changes by $\sigma(b_i x + c_i) - \sigma(b_j x + c_j) = 0$. All other terms in $f_w$ are unchanged, so $D_{(1,-1)}\Phi(w) = 0$. Again, no conditions on $\sigma$ are needed. \qedhere
\end{enumerate}
\end{solution}

\begin{solution}[ex:score-properties]
\begin{enumerate}
    \item
        Since $p(y \mid x, w) > 0$, the logarithm is well-defined and
        the chain rule gives
        \begin{equation*}
            \nabla_w \log p(y \mid x, w)
            = \frac{\nabla_w p(y \mid x, w)}{p(y \mid x, w)}.
        \end{equation*}
        The left-hand side is $s(x,y,w)$ by definition.

    \item
        Differentiate the normalisation identity
        $\int_\Y p(y \mid x, w) \, dy = 1$
        with respect to $w$:
        \begin{equation*}
            0
            = \nabla_w \int_\Y p(y \mid x, w) \, dy
            = \int_\Y \nabla_w p(y \mid x, w) \, dy.
        \end{equation*}
        By part~(a), $\nabla_w p = p \cdot s$, so
        \begin{equation*}
            0
            = \int_\Y p(y \mid x, w)\, s(x, y, w) \, dy
            = \E_{y \sim p(y \mid x, w)}\bigl[s(x, y, w)\bigr].
        \end{equation*}

    \item
        If $D_v \Psi(w) = 0$, then by the definition of
        degeneracy \eqref{eq:dist-degeneracy},
        $D_v p(y \mid x, w) = 0$ for all $x, y$.
        By part~(a), $\nabla_w p = p \cdot s$, so contracting both sides
        with $v / \|v\|$ gives $D_v p = p \cdot s_v$.
        Since $p > 0$, we conclude $s_v(x, y, w) = 0$ for all $x, y$.

    \item
        If $s_v(x, y, w) = 0$ for all $x, y$, then by part~(a),
        \begin{equation*}
            D_v p(y \mid x, w)
            = p(y \mid x, w) \cdot s_v(x, y, w)
            = 0
        \end{equation*}
        for all $x, y$ (using $p > 0$). That is,
        $D_v \Psi(w) = 0$. \qedhere
\end{enumerate}
\end{solution}

\begin{solution}[ex:fim-degeneracy]
\begin{enumerate}
    \item
        Since $v$ is a unit vector,
        $s_v = D_v \log p = \frac{v}{\|v\|} \cdot s = v \cdot s$.
        From the definition of $I(w)$ in \eqref{eq:fim},
        \begin{align*}
            v^\top I(w)\, v
            &= v^\top
               \E_{x,y}\!\bigl[s\, s^\top\bigr]\, v
            = \E_{x,y}\!\bigl[(v \cdot s)^2\bigr]
            = \E_{x,y}\!\bigl[s_v^2\bigr]
        \\  &= \E_{x \sim q(x)}\,
               \E_{y \sim p(y \mid x, w)}\!\bigl[
                   s_v(x,y,w)^2
               \bigr].
        \end{align*}

    \item
        If $D_v \Psi(w) = 0$ for some non-zero $v$, then by
        \cref{ex:score-properties:b}, $s_{v}(x,y,w) = 0$ for all $x, y$.
        Let $\hat{v} = v / \|v\|$. Then $s_{\hat{v}} = s_v$ (the
        directional derivative is invariant to the magnitude of $v$), so
        by part~(a),
        $\hat{v}^\top I(w)\, \hat{v} = \E[s_{\hat{v}}^2] = 0$.
        A matrix $M$ is positive definite only if $u^\top M u > 0$ for
        all non-zero $u$. Since $\hat{v}$ is a non-zero vector with
        $\hat{v}^\top I(w)\, \hat{v} = 0$, the matrix $I(w)$ is not
        positive definite.

    \item
        If $I(w)$ is not positive definite, then since $I(w)$ is
        positive semidefinite, there exists a non-zero vector $\hat{v}$
        such that $\hat{v}^\top I(w)\, \hat{v} = 0$. By part~(a) (taking
        $\hat{v}$ to be a unit vector without loss of generality),
        $\E[s_{\hat{v}}^2] = 0$. Since $s_{\hat{v}}^2 \geq 0$ and
        $q(x) > 0$ and $p(y \mid x, w) > 0$, this implies
        $s_{\hat{v}}(x,y,w) = 0$ for all $x \in \X$ and $y \in \Y$.
        By \cref{ex:score-properties:c}, $D_{\hat{v}} \Psi(w) = 0$.

    \item
        Parts~(b) and~(c) show that $I(w)$ is positive definite if and
        only if $\Psi$ is non-degenerate at $w$, or equivalently, $I(w)$
        is singular if and only if $\Psi$ is degenerate at $w$.

        For the kernel characterisation: clearly $0 \in \ker I(w)$.
        For non-zero $v$, $v \in \ker I(w)$ means $I(w)\,v = 0$, which
        (since $I(w)$ is positive semidefinite) is equivalent to
        $v^\top I(w)\, v = 0$. Normalising to $\hat{v} = v/\|v\|$,
        part~(a) gives $\E[s_{\hat{v}}^2] = 0$, which as in part~(c)
        implies $s_{\hat{v}} = 0$ everywhere. By
        \cref{ex:score-properties:c}, $D_{\hat{v}} \Psi(w) = 0$, and
        hence $D_v \Psi(w) = 0$. Conversely, if $D_v \Psi(w) = 0$ for
        non-zero $v$, then \cref{ex:score-properties:b} gives $s_v = 0$,
        so $v^\top I(w)\, v = \|v\|^2 \E[s_v^2] = 0$, hence
        $I(w)\, v = 0$. Therefore
        $\ker I(w)
        = \{v \in \R^d : D_v \Psi(w) = 0\}$. \qedhere
\end{enumerate}
\end{solution}

\begin{solution}[ex:loss-basic]
\begin{enumerate}
    \item
        Since $L_{x,y}(w)$ depends on $w$ only through $\Phi(w)$, we can write
        $L_{x,y}(w) = \ell_{x,y}(\Phi(w))$ for some functional $\ell_{x,y}$.
        If $D_v\Phi(w) = 0$, then
            $\Phi(w + \epsilon v) = \Phi(w) + O(\epsilon^2)$
        as $\epsilon \to 0$, and therefore
            $L_{x,y}(w + \epsilon v)
            = \ell_{x,y}(\Phi(w) + O(\epsilon^2))
            = L_{x,y}(w) + O(\epsilon^2)$.
        Taking expectations over $(x,y) \sim q$,
        \begin{equation*}
            L(w + \epsilon v) = L(w) + O(\epsilon^2),
        \end{equation*}
        so $D_v L(w)
        = \lim_{\epsilon \to 0}
            \frac{L(w + \epsilon v) - L(w)}{\epsilon\|v\|}
        = 0$.

    \item
        If $\nabla L(w) = 0$, then $D_v L(w) = \frac{v}{\|v\|} \cdot
        \nabla L(w) = 0$ for every nonzero $v$.

        If $\nabla L(w) \neq 0$, the orthogonal complement
        $\nabla L(w)^\perp$ has dimension $d - 1 \geq 1$ (since $d > 1$).
        Pick any nonzero $v \in \nabla L(w)^\perp$. Then
        $D_v L(w) = \frac{v}{\|v\|} \cdot \nabla L(w) = 0$. \qedhere
\end{enumerate}
\end{solution}

\begin{solution}[ex:loss-landscape-degeneracy]
\begin{enumerate}
    \item
        Since $a^{2k} \geq 0$ and $b^{2l} \geq 0$ for all $a, b \in \R$
        (using the convention $0^0 = 1$), we have $L_{k,l}(a,b) \geq 0$ for
        $k, l \geq 1$, $L_{k,l}(a,b) \geq 1$ if $k = 0$ or $l = 0$, and
        $L_{k,l}(a,b) = 2$ if $k = l = 0$.
        At the origin,
        $L_{k,l}(0,0) = 0^{2k} + 0^{2l}$,
        which equals $0$ if $k \geq 1$ and $l \geq 1$, equals $1$ if exactly
        one of $k, l$ is $0$, and equals $2$ if $k = l = 0$.
        In each case this matches the lower bound, so the origin is a global
        minimum.

        The origin is the \emph{unique} global minimum if and only if
        $k \geq 1$ and $l \geq 1$: then $L_{k,l}(a,b) = 0$ requires both
        $a^{2k} = 0$ and $b^{2l} = 0$, forcing $a = b = 0$.
        If $k = 0$, then $L_{0,l}(a,0) = 1$ for all $a$, so the entire
        $a$-axis achieves the minimum. Similarly if $l = 0$. If $k = l = 0$,
        then all parameters achieve the minimum.

    \item
        For $k = l = 1$: $L_{1,1}(a,b) = a^2 + b^2$. The Hessian at the
        origin is
        \begin{equation*}
            H(0,0) = \begin{pmatrix} 2 & 0 \\ 0 & 2 \end{pmatrix},
        \end{equation*}
        which is positive definite. The origin is a non-degenerate minimum.

    \item
        For $k > 1$ and $l \geq 1$:
        $\frac{\partial^2 L_{k,l}}{\partial a^2}
        = 2k(2k{-}1)\, a^{2k-2}$.
        At the origin, $a^{2k-2} = 0$ since $2k - 2 \geq 2$, so this entry
        vanishes. The Hessian at the origin is therefore
        \begin{equation*}
            H(0,0) = \begin{pmatrix}
                0 & 0 \\
                0 & 2l(2l{-}1) \cdot 0^{2l-2}
            \end{pmatrix},
        \end{equation*}
        which has a zero eigenvalue (from the first diagonal entry) regardless
        of the value of the second. The origin is a degenerate minimum.

    \item
        For $k = 0$ and $l \geq 1$: $L_{0,l}(a,b) = 1 + b^{2l}$.
        Since $L_{0,l}$ is independent of $a$, both $\frac{\partial L}{\partial
        a}$ and $\frac{\partial^2 L}{\partial a^2}$ vanish identically. The
        Hessian at the origin has a zero first diagonal entry, so it is
        singular. The origin is a degenerate minimum. \qedhere
\end{enumerate}
\end{solution}

\begin{solution}[ex:fim-hessian]
\begin{enumerate}
    \item
        From \cref{ex:score-properties}, $\E_{y \sim p(y \mid x, w)}[s_j(x,y,w)]
        = 0$ for each component $j$ of the score, where
        $s_j = \frac{\partial}{\partial w_j} \log p(y \mid x, w)$.
        Differentiating with respect to $w_k$ and exchanging the derivative
        with the integral:
        \begin{align*}
            0
            &= \frac{\partial}{\partial w_k}
                \int_\Y s_j(x,y,w)\, p(y \mid x, w)\, dy
        \\  &= \int_\Y
                \frac{\partial s_j}{\partial w_k}\, p\, dy
                + \int_\Y s_j\, \frac{\partial p}{\partial w_k}\, dy.
                \tag{product rule}
        \end{align*}
        Using the score identity
        $\frac{\partial p}{\partial w_k} = p \cdot s_k$
        \eqref{eq:score-identity}, the second integral becomes
        $\int_\Y s_j\, s_k\, p\, dy = \E_y[s_j\, s_k]$. Therefore
        \begin{equation*}
            \E_{y \sim p(y \mid x, w)}\!\left[
                \frac{\partial^2 \log p}{\partial w_j \partial w_k}
            \right]
            = -\E_{y \sim p(y \mid x, w)}[s_j\, s_k].
        \end{equation*}
        Assembling all components into a matrix gives
        \begin{equation*}
            -\E_{y \sim p(y \mid x, w)}\!\bigl[
                \nabla_w^2 \log p(y \mid x, w)
            \bigr]
            =
            \E_{y \sim p(y \mid x, w)}\!\bigl[
                s(x,y,w)\, s(x,y,w)^\top
            \bigr].
        \end{equation*}

    \item
        The population negative log-likelihood is
        $L(w) = -\E_{x \sim q(x)}\, \E_{y \sim p(y \mid x, w_0)}
            [\log p(y \mid x, w)]$.
        Taking the Hessian with respect to $w$ (exchanging differentiation and
        integration):
        \begin{equation*}
            H(w)
            = -\E_{x \sim q(x)}\,
              \E_{y \sim p(y \mid x, w_0)}\!\bigl[
                \nabla_w^2 \log p(y \mid x, w)
              \bigr].
        \end{equation*}
        At $w = w_0$, the inner expectation is over
        $y \sim p(y \mid x, w_0)$, which matches the distribution in the
        Bartlett identity. Applying part~(a):
        \begin{align*}
            H(w_0)
            &= \E_{x \sim q(x)}\,
              \E_{y \sim p(y \mid x, w_0)}\!\bigl[
                s(x,y,w_0)\, s(x,y,w_0)^\top
              \bigr]
        \\  &= I(w_0).
            \tag{by \cref{def:fim}}
        \end{align*}

    \item
        By \cref{ex:fim-degeneracy},
        $\ker I(w_0)
        = \{v \in \R^d : D_v \Psi(w_0) = 0\}$.
        Since $H(w_0) = I(w_0)$, if $\Psi$ is degenerate at $w_0$ in
        direction $v$, then $v \in \ker I(w_0) = \ker H(w_0)$, so
        $H(w_0)\, v = 0$.

        Contrapositively: if $H(w_0)$ is positive definite, then
        $\ker H(w_0) = \{0\}$, so $\ker I(w_0) = \{0\}$, and $\Psi$ is
        non-degenerate at $w_0$. \qedhere
\end{enumerate}
\end{solution}

\begin{solution}[ex:loss-degeneracy-examples]
\begin{enumerate}
    \item
        The identity parameter--function map $\Phi(a,b) = (a,b)$ has Jacobian
        equal to the $2 \times 2$ identity matrix everywhere, so it is
        non-degenerate at every point.
        \Cref{ex:lld:1} shows that $L_{1,1}$ has a non-degenerate minimum:
        absence of parameter--function map degeneracy is consistent with
        absence of loss landscape degeneracy.
        \Cref{ex:lld:2,ex:lld:3} show that $L_{2,1}$, $L_{0,1}$, etc., have
        degenerate minima despite the parameter--function map being
        non-degenerate. This demonstrates that loss landscape degeneracy does
        not imply parameter--function map degeneracy.

    \item
        $L_\circ(a,b) = (a^2 + b^2 - 1)^2 \geq 0$, with equality if and only
        if $a^2 + b^2 = 1$. So every point on the unit circle is a global
        minimum. Consider $(1, 0)$. The gradient is
        \begin{equation*}
            \nabla L_\circ
            = \bigl(4a(a^2 + b^2 - 1),\; 4b(a^2 + b^2 - 1)\bigr),
        \end{equation*}
        which vanishes at $(1, 0)$. The Hessian entries at $(1, 0)$ are:
        \begin{align*}
            \frac{\partial^2 L_\circ}{\partial a^2}
            &= 4(3a^2 + b^2 - 1) = 8,
        &\qquad
            \frac{\partial^2 L_\circ}{\partial b^2}
            &= 4(a^2 + 3b^2 - 1) = 0,
        \\
            \frac{\partial^2 L_\circ}{\partial a\, \partial b}
            &= 8ab = 0.
        \end{align*}
        So $H(1,0) = \mathrm{diag}(8, 0)$, which is singular. Therefore, $(1,0)$
        is a degenerate minimum.

        This provides another example of loss landscape degeneracy without
        parameter--function map degeneracy: the identity parameter--function
        map is non-degenerate, but the circle of global minima creates a flat
        tangential direction.

    \item
    \begin{enumerate}
        \item
            The map $\Phi(\alpha, \beta) = (\alpha^3, \beta^3)$ is
            the product of two copies of the cubic parametrisation from
            \cref{ex:cubic-degeneracy}. In each component, the derivative
            $\frac{d}{d\alpha}(\alpha^3) = 3\alpha^2$ vanishes at
            $\alpha = 0$. Therefore both coordinate directions are degenerate
            at the origin, and hence every direction is degenerate at the
            origin (since the Jacobian $\mathrm{diag}(3\alpha^2, 3\beta^2)$
            is the zero matrix there).

        \item
            In the new parametrisation, the loss becomes
            $L_{k,l}(\alpha,\beta) = (\alpha^3)^{2k} + (\beta^3)^{2l}
            = \alpha^{6k} + \beta^{6l}$.

            If $k \geq 1$ and $l \geq 1$: the exponents $6k$ and $6l$ are
            both at least $6$, so
            $\frac{\partial^2 L}{\partial \alpha^2}\big|_0
            = 6k(6k{-}1) \cdot 0^{6k-2} = 0$
            (since $6k - 2 \geq 4$),
            and similarly for $\beta$. The Hessian at the origin is the zero
            matrix, which is singular.

            If $k = 0$: $L_{0,l}(\alpha,\beta) = 1 + \beta^{6l}$, which is
            independent of $\alpha$, so the Hessian has a zero first diagonal
            entry. Similarly if $l = 0$. In all cases, the origin is a
            degenerate minimum.

        \item
            Under the identity parameter--function map, $L_{1,1}(a,b) = a^2 +
            b^2$ had a non-degenerate minimum at the origin (\cref{ex:lld:1}).
            Under the cubic parameter--function map, the same loss became
            $\alpha^6 + \beta^6$, which has a degenerate minimum at the origin.
            This shows that introducing degeneracy into the parameter--function
            map can convert a non-degenerate minimum into a degenerate one.
            
            For other $k,l$, $L_{k,l}$ was already degenerate at the origin
            under the non-degenerate parameterisation, the cubic function only
            makes it \emph{more} degenerate (in terms of the order to which the
            loss vanishes in the degenerate direction(s); we will formally
            quantify this in \cref{sec:llc}).
    \end{enumerate}

    \item
    \begin{enumerate}
        \item
            The gradient is
            $\nabla L = (ab - 1)(b,\, a)$.
            At $(0, 0)$: $\nabla L = (0 - 1)(0, 0) = (0, 0)$.
            So $(0, 0)$ is a critical point.

            The Jacobian of $\Phi(a,b) = ab$ is $(b, a)$, which at
            $(0, 0)$ is $(0, 0)$. Every direction $v = (v_1, v_2)$ gives
            $J v = b v_1 + a v_2 = 0$, so $\Phi$ is degenerate in every
            direction at $(0, 0)$.

        \item
            Computing second partial derivatives:
            \begin{align*}
                \frac{\partial^2 L}{\partial a^2} &= b^2,
            &\qquad
                \frac{\partial^2 L}{\partial b^2} &= a^2,
            &\qquad
                \frac{\partial^2 L}{\partial a\, \partial b}
                &= 2ab - 1.
            \end{align*}
            At $(0, 0)$:
            \begin{equation*}
                H(0,0) = \begin{pmatrix} 0 & -1 \\ -1 & 0 \end{pmatrix},
            \end{equation*}
            which has determinant $-1 \neq 0$ (eigenvalues $\pm 1$).
            The Hessian is nonsingular.

        \item
            Despite full degeneracy of the parameter--function map at $(0,0)$,
            the Hessian is nonsingular (in fact indefinite, so $(0,0)$ is a
            saddle point, not a local minimum). This shows that
            parameter--function map degeneracy does not in general imply loss
            landscape degeneracy.

            This does not contradict \cref{ex:fim-hessian}, which establishes
            $H(w_0) = I(w_0)$ at the \emph{true parameter} $w_0$, meaning
            $\Phi(w_0)$ must equal the target. Here, $\Phi(0,0) = 0 \neq 1$,
            so $(0,0)$ is not the true parameter, and the identity does not
            apply. \qedhere
    \end{enumerate}
\end{enumerate}
\end{solution}



% \clearpage
% \subsection[Solutions to exercises from Section~3]{Solutions to exercises from \cref{sec:llc}}
% \addtocontents{loe}{\solutionlistskip}

\begin{solution}[ex:llc-regular-case]
    Because $w^*$ is a local minimum, the gradient $\nabla L(w^*)$ is zero. By Taylor's theorem, we can write the loss function near $w^*$ exactly as:
    $$L(w) - L(w^*) = \frac{1}{2} (w - w^*)^\top H (w - w^*) + o(\|w - w^*\|^2)$$
    where $H = \nabla^2 L(w^*)$ is the strictly positive definite Hessian matrix.
    
    The sublevel set is defined by the condition $L(w) - L(w^*) < \epsilon$. As $\epsilon \to 0$, the neighborhood of parameters that satisfy this condition shrinks toward $w^*$. In this limit, the higher-order error term $o(\|w - w^*\|^2)$ becomes negligible, and the boundary of the sublevel set is governed entirely by the quadratic form:
    $$\frac{1}{2} (w - w^*)^\top H (w - w^*) < \epsilon.$$
    
    We can rewrite this inequality into the standard form of an ellipsoid, $x^\top A x < 1$, by dividing both sides by $\epsilon$:
    $$(w - w^*)^\top \left( \frac{1}{2\epsilon} H \right) (w - w^*) < 1.$$
    
    Using the hint, the volume of a $d$-dimensional ellipsoid defined by a matrix $A$ is proportional to $(\det A)^{-1/2}$. Here, our matrix is $A = \frac{1}{2\epsilon}H$. 
    
    Therefore,
    $$\det\left( \frac{1}{2\epsilon} H \right)^{-1/2} = \left( \left(\frac{1}{2\epsilon}\right)^d \det H \right)^{-1/2} = \epsilon^{d/2} \underbrace{\left( \frac{\det H}{2^d} \right)^{-1/2}}_{:=c} $$
    
    Because the true loss function deviates from this pure quadratic form only by an error of $o(\|w - w^*\|^2)$, the volume of the true sublevel set deviates from the ellipsoid's volume only by a correspondingly higher-order error term. Thus, we conclude:
    $$V(\epsilon) = c \epsilon^{d/2} + o(\epsilon^{d/2}) \quad \text{as } \epsilon \to 0.$$
\end{solution}

\begin{solution}[ex:quad-valley-scaling]
\begin{enumerate}
    \item The loss is $L(x_1, x_2) = x_1^2$, which is a parabolic trough: a parabola in the $x_1$ direction and constant along $x_2$. The zero set is the line $x_1 = 0$.
    \item $\W_0 = \{0\}^{d'} \times [-R,R]^{d-d'}$, which is a $(d - d')$-dimensional affine subspace of $\W$. Hence $\dim(\W_0) = d - d'$.
    \item Since $L(w) = \sum_{i=1}^{d'} x_i^2$, the partial derivatives are
    \[
        \frac{\partial L}{\partial x_i} = \begin{cases} 2x_i & i \leq d', \\ 0 & i > d', \end{cases}
        \qquad
        \frac{\partial^2 L}{\partial x_i \partial x_j} = \begin{cases} 2 & i = j \leq d', \\ 0 & \text{otherwise}. \end{cases}
    \]
    So the Hessian is the block diagonal matrix:
    \[
        \nabla^2 L = \begin{pmatrix} 2I_{d'} & 0 \\ 0 & 0_{d-d'} \end{pmatrix},
    \]
    which has rank $d'$. In particular it is positive semi-definite but not positive definite whenever $d' < d$.
    \item The minimum of $L$ is $0$, achieved on $\W_0$. For any $w^* \in \W_0$ and $\epsilon > 0$, the sublevel set is
    \[
        B(w^*, \epsilon) = \left\{ (x_1,\dots,x_d) \in [-R,R]^d \;\middle|\; \sum_{i=1}^{d'} x_i^2 < \epsilon \right\} = \left\{x \in \R^{d'} : \sum_{i=1}^{d'} x_i^2 < \epsilon \right\} \times [-R,R]^{d-d'}
    \]
    for $R$ sufficiently large (specifically $\sqrt{\epsilon} < R$). The condition $\sum_{i=1}^{d'} x_i^2 < \epsilon$ can be rewritten as $\mathbf{x}^\top A \mathbf{x} < 1$ where $A = \epsilon^{-1} I_{d'}$. By the ellipsoid volume formula from \cref{ex:llc-regular-case}, the volume of this region is proportional to $(\det A)^{-1/2} = (\epsilon^{-d'})^{-1/2} = \epsilon^{d'/2}$. Multiplying by the volume $(2R)^{d-d'}$ of the remaining coordinates gives $V(\epsilon) \propto \epsilon^{d'/2}$. \qedhere
\end{enumerate}
\end{solution}

\begin{solution}[ex:solving-for-lambda]
\begin{enumerate}
    \item
    From \cref{def:llc}, as $\epsilon \to 0$,
    \begin{equation*}
        V(\epsilon)
        = c\,\epsilon^{\lambda}(-\log\epsilon)^{m-1}\bigl(1 + o(1)\bigr).
    \end{equation*}
    Substituting $\epsilon \mapsto a\epsilon$,
    \begin{equation*}
        V(a\epsilon)
        = c\,(a\epsilon)^{\lambda}
            (-\log(a\epsilon))^{m-1}\bigl(1 + o(1)\bigr).
    \end{equation*}
    Taking the ratio,
    \begin{equation*}
        \frac{V(a\epsilon)}{V(\epsilon)}
        = a^{\lambda}
          \left(\frac{-\log(a\epsilon)}{-\log\epsilon}\right)^{\!m-1}
          \bigl(1 + o(1)\bigr).
    \end{equation*}
    As $\epsilon \to 0$, $-\log\epsilon \to \infty$, and
    \begin{equation*}
        \frac{-\log(a\epsilon)}{-\log\epsilon}
        = \frac{-\log a - \log\epsilon}{-\log\epsilon}
        = 1 + \frac{\log a}{-\log\epsilon}
        \;\xrightarrow{\epsilon\to 0}\; 1.
    \end{equation*}
    Therefore $V(a\epsilon)/V(\epsilon) \to a^{\lambda}$. Taking
    logarithms and dividing by $\log a > 0$ gives the result.
    \item From \cref{def:llc},
    $V(\epsilon) = c\,\epsilon^{\lambda(w^*)}
        (-\log\epsilon)^{\mu(w^*)-1}(1 + o(1))$,
    so
    \begin{equation*}
        \log V(\epsilon)
        = \lambda(w^*)\log\epsilon
          + (\mu(w^*)-1)\log(-\log\epsilon)
          + \log c + o(1).
    \end{equation*}
    Dividing by $\log\epsilon$
    \begin{equation*}
        \frac{\log V(\epsilon)}{\log\epsilon}
        = \lambda(w^*)
        + \frac{(\mu(w^*)-1)\log(-\log\epsilon)}{\log\epsilon}
        + \frac{\log c + o(1)}{\log\epsilon}.
    \end{equation*}

    It remains to prove that the final two terms in this sum vanish as $\epsilon \to 0$. The final term tends to $0$ since $\log \epsilon \to -\infty$ and the numerator is bounded as $\epsilon \to 0$. Now we look at the middle term which is proportional to
    $$\lim\limits_{\epsilon\to 0} \frac{\log (-\log \epsilon)}{\log \epsilon}=\lim\limits_{x\to \infty} \frac{\log (x)}{-x} = 0,$$ completing the proof. \qedhere

\end{enumerate}
\end{solution}

\begin{solution}[ex:calculating-llc]
\begin{enumerate}
    \item
        $V(\epsilon) = \operatorname{Vol}(\{x : x^2 < \epsilon\}) = 2\sqrt{\epsilon}$, so $V(\epsilon) \propto \epsilon^{1/2}$ and $\lambda = 1/2$. The Hessian is $L''(0) = 2$, which has rank $1$ (full rank).
    \item
        $V(\epsilon) = \operatorname{Vol}(\{x : x^{2m} < \epsilon\}) = 2\epsilon^{1/2m}$, so $\lambda = 1/2m$. For $m \geq 2$, the first $2m - 1$ derivatives vanish at $x = 0$, so the Hessian is $L''(0) = 0$, which has rank $0$. Since $\lambda = 1/2m < 1/2 = d/2$, the singularity makes the model effectively simpler than its parameter count suggests.
    \item
        The sublevel set is $\{(x,y) : x^4 + y^4 < \epsilon\}$. Substituting $x = \epsilon^{1/4}s$, $y = \epsilon^{1/4}t$:
        \[
            V(\epsilon) = \epsilon^{1/2} \int\!\!\int_{\{s^4 + t^4 < 1\}} ds\, dt = c\,\epsilon^{1/2},
        \]
        so $\lambda = 1/2$. The Hessian at the origin is $\nabla^2L(0,0) = \operatorname{diag}(12x^2, 12y^2)\big|_{(0,0)} = 0$, which has rank $0$.
    \item
        The sublevel set is $\{(x,y) : x^2 + y^4 < \epsilon\}$. For each fixed $y$ with $|y| < \epsilon^{1/4}$, $x$ ranges over $|x| < \sqrt{\epsilon - y^4}$. Substituting $y = \epsilon^{1/4}t$:
        \[
            V(\epsilon) = \int_{-\epsilon^{1/4}}^{\epsilon^{1/4}} 2\sqrt{\epsilon - y^4}\, dy = 2\epsilon^{1/4} \int_{-1}^{1} \sqrt{\epsilon - \epsilon t^4}\, dt = 2\epsilon^{3/4} \int_{-1}^{1} \sqrt{1 - t^4}\, dt = c\,\epsilon^{3/4},
        \]
        so $\lambda = 3/4$. The Hessian is $\nabla^2L(0,0) = \operatorname{diag}(2, 12y^2)\big|_{(0,0)} = \operatorname{diag}(2, 0)$, which has rank $1$.
\end{enumerate}
\end{solution}


\begin{solution}[ex:lambda-upper-bound]
\begin{enumerate}
    \item 
        Since $w^*$ is a local minimum of $L$, we have $\nabla L(w^*) = 0$. By Taylor's theorem with Lagrange remainder, for any $w \in B(w^*)$ there exists $\tilde{w}$ on the line segment between $w^*$ and $w$ such that
        \begin{align*}
            L(w) - L(w^*) 
            &= \nabla L(w^*)^\top (w - w^*) + \frac{1}{2}(\tilde w - w^*)^\top \nabla^2 L(\tilde{w})(\tilde w - w^*)
            \\&= \frac{1}{2}(\tilde w - w^*)^\top \nabla^2 L(\tilde{w})(\tilde w - w^*).
        \end{align*}
        Since $\tilde{w} \in B(w^*)$, the eigenvalues of $\nabla^2 L(\tilde{w})$ are bounded above by $\Lambda$, so
        \begin{equation*}
            (\tilde w - w^*)^\top \nabla^2 L(\tilde{w})(\tilde w - w^*)
            \leq \Lambda \|\tilde w - w^*\|^2
            \leq \Lambda \|w - w^*\|^2.
        \end{equation*}
        Combining, $L(w) - L(w^*) \leq \frac{1}{2}\Lambda \|w - w^*\|^2$.

    \item 
        Set $r = \sqrt{2\epsilon/\Lambda}$. Set $\epsilon$ sufficiently small that $B_r(w^*) \subseteq B(w^*)$.
        If $w \in B_r(w^*)$, then $\|w - w^*\| < r$, so by part (a),
        \begin{equation*}
            L(w) - L(w^*) \leq \frac{1}{2}\Lambda \|w - w^*\|^2 < \frac{1}{2}\Lambda \cdot \frac{2\epsilon}{\Lambda} = \epsilon.
        \end{equation*}
        Hence $w \in B(w^*, \epsilon)$, and therefore $B_r(w^*) \subseteq B(w^*, \epsilon)$.

    \item 
        Since $B_r(w^*) \subseteq B(w^*, \epsilon)$, 
        \begin{equation*}
            V(\epsilon) 
            = \operatorname{Vol}(B(w^*, \epsilon)) 
            \geq \operatorname{Vol}(B_r(w^*)) 
            = c \, r^d 
            = c \left(\frac{2\epsilon}{\Lambda}\right)^{d/2} 
            = C \, \epsilon^{d/2},
        \end{equation*}
        for some $c >0$ and $C = c (2/\Lambda)^{d/2} > 0$. This holds for all sufficiently small $\epsilon > 0$ (i.e.\ small enough that $B_r(w^*) \subseteq B(w^*)$).

    \item 
        From part (c), $V(\epsilon) \geq C\epsilon^{d/2}$ for sufficiently small $\epsilon$ (w.l.o.g.\ assume $\epsilon < 1$). Taking logarithms,
        \begin{equation*}
            \log V(\epsilon) \geq \log C + \frac{d}{2}\log \epsilon.
        \end{equation*}
        Dividing by $\log \epsilon < 0$ (since $\epsilon < 1$) reverses the inequality:
        \begin{equation*}
            \frac{\log V(\epsilon)}{\log \epsilon} \leq \frac{\log C}{\log \epsilon} + \frac{d}{2}.
        \end{equation*}
        Since $\log C / \log \epsilon \to 0$ as $\epsilon \to 0^+$, taking the limit and using Equation \eqref{eq:solve-LLC} gives
        \begin{equation*}
            \lambda(w^*) = \lim_{\epsilon \to 0^+} \frac{\log V(\epsilon)}{\log \epsilon} \leq \frac{d}{2}. \qedhere
        \end{equation*}
\end{enumerate}
\end{solution}

\begin{solution}[ex:LLC-hessian-rank]

    Let $L_1(x,y) = x^2+y^4$ with $w_1 = (0,0)$. Then $\nabla^2L_1(0,0) = \operatorname{diag}(2, 0)$, so $\operatorname{rank}\nabla^2L_1(0,0) = 1$. By \cref{ex:calculating-llc}, $\lambda(w_1) = \frac{3}{4}$.

    Let $L_2(x,y) = x^2$ with $w_2 = (0,0)$. Then $\nabla^2L_2(0,0) = \operatorname{diag}(2, 0)$, so $\operatorname{rank}\nabla^2L_2(0,0) = 1$. The minimiser locus is $\{x = 0\}$, which is the setting of \cref{ex:quad-valley-scaling} with $d' = 1$, giving $\lambda(w_2) = \frac{1}{2}$.

    Both Hessians have rank $1$, but $\lambda(w_1) = \frac{3}{4} \neq \frac{1}{2} = \lambda(w_2)$. The difference is that the flat direction in $L_2$ is \emph{exactly} flat (the minimisers form a line), while in $L_1$ it is only flat to second order but quartic beyond that. The LLC detects this distinction; the Hessian rank does not.
\end{solution}

\begin{solution}[ex:cubic-parameterisation]

\begin{enumerate}
    \item
        The cubically-parameterised loss $L(\mu) = \tfrac{1}{2}(\mu^3 -
        \mu_0^3)^2$ vanishes if and only if $\mu^3 = \mu_0^3$, i.e.\
        $\mu = \mu_0$ (since $t \mapsto t^3$ is a bijection on $\R$).
        The ordinary quadratic loss $\ell(\theta) =
        \tfrac{1}{2}(\theta - \theta_0)^2$ vanishes if and only if
        $\theta = \theta_0$. Since for every $\theta_0 \in \R$ we can
        write $\theta_0 = \mu_0^3$ for a unique $\mu_0 = \theta_0^{1/3}$,
        both losses achieve their minimum value of zero at exactly the
        same set of ``targets'' $\theta_0 \in \R$, so the set of
        achievable minima is the same.

    \item
        Computing derivatives:
        \begin{align*}
            L'(\mu) &= 3\mu^2(\mu^3 - \mu_0^3),
            \\
            L''(\mu) &= 6\mu(\mu^3 - \mu_0^3) + 9\mu^4.
        \end{align*}
        At the minimum $\mu = \mu_0$:
        $L''(\mu_0) = 6\mu_0 \cdot 0 + 9\mu_0^4 = 9\mu_0^4$.
        \begin{itemize}
            \item If $\mu_0 = 0$: $L''(0) = 0$, so the Hessian
                vanishes and the loss is degenerate at its minimum.
            \item If $\mu_0 \neq 0$: $L''(\mu_0) = 9\mu_0^4 > 0$,
                so the loss is non-degenerate.
        \end{itemize}
    
    \item
        We compute $V(\epsilon) = \mathrm{vol}(\{\mu \in \R : L(\mu) < \epsilon\})$ in each case.

\medskip
\noindent\textit{Case $\mu_0 = 0$.}
When $\mu_0 = 0$ the loss simplifies to $L(\mu) = \tfrac{1}{2}\mu^6$, so
\[
    V(\epsilon)
    = \int_{\frac{1}{2}\mu^6 < \epsilon} d\mu
    = \int_{|\mu| < (2\epsilon)^{1/6}} d\mu
    = 2(2\epsilon)^{1/6}
    = 2^{7/6}\,\epsilon^{1/6}.
\]

\medskip
\noindent\textit{Case $\mu_0 \neq 0$.}
The sublevel set condition $L(\mu) < \epsilon$ is $\tfrac{1}{2}(\mu^3 - \mu_0^3)^2 < \epsilon$, i.e.\
\[
    \mu_0^3 - \sqrt{2\epsilon} \;<\; \mu^3 \;<\; \mu_0^3 + \sqrt{2\epsilon}.
\]
Since the cube root is monotone on $\R$, this is equivalent to
\[
    \bigl(\mu_0^3 - \sqrt{2\epsilon}\bigr)^{1/3}
    \;<\; \mu \;<\;
    \bigl(\mu_0^3 + \sqrt{2\epsilon}\bigr)^{1/3},
\]
and therefore
\[
    V(\epsilon)
    = \bigl(\mu_0^3 + \sqrt{2\epsilon}\bigr)^{1/3}
    - \bigl(\mu_0^3 - \sqrt{2\epsilon}\bigr)^{1/3}.
\]

    \item
        \textbf{Case $\mu_0 = 0$:}
        $V(\epsilon) \propto \epsilon^{1/6}$, so
        $\lambda = \tfrac{1}{6}$.

        \textbf{Case $\mu_0 \neq 0$:}
        For small $\epsilon$ we expand by factoring out $\mu_0$:
\[
    V(\epsilon)
    = \mu_0 \left[
        \left(1 + \frac{\sqrt{2\epsilon}}{\mu_0^3}\right)^{\!1/3}
        - \left(1 - \frac{\sqrt{2\epsilon}}{\mu_0^3}\right)^{\!1/3}
    \right].
\]
Applying the binomial expansion $(1+x)^{1/3} = 1 + \tfrac{1}{3}x + o(x)$ to each term with $x = \pm\frac{\sqrt{2\epsilon}}{\mu_0^3}$, the constant terms cancel and the linear terms add:
\[
    V(\epsilon)
    = \mu_0 \left[
        \frac{1}{3}\frac{\sqrt{2\epsilon}}{\mu_0^3}
        + \frac{1}{3}\frac{\sqrt{2\epsilon}}{\mu_0^3}
        + o(\epsilon^{1/2})
    \right]
    = \frac{2\sqrt{2}}{3\mu_0^2}\,\epsilon^{1/2} + o(\epsilon^{1/2}).
\]
        In the non-degenerate case $\mu_0 \neq 0$, we recover the regular
        value $\lambda = d/2 = 1/2$ as expected.
        In the degenerate case $\mu_0 = 0$, $\lambda = 1/6 < 1/2$: the
        singularity at $\mu = 0$ (where the loss vanishes to sixth order
        rather than second order) makes the model effectively simpler than
        its parameter count suggests.

    \item From part (c), for fixed $\epsilon > 0$ we plot
\[
    V(\epsilon;\, \mu_0)
    = \bigl(\mu_0^3 + \sqrt{2\epsilon}\bigr)^{1/3}
    - \bigl(\mu_0^3 - \sqrt{2\epsilon}\bigr)^{1/3}
\]

The function is a smooth, symmetric bump centred at $\mu_0 = 0$ with no discontinuity. Making $\epsilon$ smaller narrows the bump and lowers the peak; making $\epsilon$ larger widens it. Despite the learning coefficient jumping discontinuously from $1/2$ to $1/6$ at $\mu_0 = 0$ exactly, the finite-$\epsilon$ volume transitions smoothly, with the singularity's influence confined to a neighbourhood that shrinks as $\epsilon \to 0$.
\end{enumerate}
\end{solution}

\begin{solution}[ex:rlct-computation]
The loss $L(w) = w_1^{2k_1}\cdots w_d^{2k_d}$ is already in monomial
form, so the identity map $g = \mathrm{id}$ serves as a
(trivial) resolution of singularities with $dw = 1 \cdot du$,
i.e.\ $h_j = 0$ for all $j$.
Substituting into \eqref{eq:rlct} gives
\[
    \lambda(0)
    \;=\; \min_{j:\, k_j > 0} \frac{0 + 1}{2k_j}
    \;=\; \min_{j:\, k_j > 0} \frac{1}{2k_j}\,. \qedhere
\]
\end{solution}

\begin{solution}[ex:holder-llc]
    Under the loss pseudo-metric $\rho(w,u) = |L(w) - L(u)|$, the
    $\epsilon$-ball centred at $w^*$ (within $B(w^*)$) is
    \begin{equation*}
        B_\epsilon(w^*)
        = \{w \in B(w^*) : |L(w) - L(w^*)| < \epsilon\}.
    \end{equation*}
    Since $w^*$ is a local minimum and $L(w) \geq L(w^*)$ on $B(w^*)$,
    this simplifies to
    $B_\epsilon(w^*) = \{w \in B(w^*) : L(w) - L(w^*) < \epsilon\}
    = B(w^*,\epsilon)$, the sublevel set from \eqref{eq:sublevel-set}. In
    particular, $\nu(B_\epsilon(w^*)) = V(\epsilon)$ where $\nu$ is
    Lebesgue measure.

    The H\"older exponent is therefore
    \begin{equation*}
        \alpha(w^*)
        = \lim_{\epsilon\to 0}
            \frac{\log \nu(B_\epsilon(w^*))}{\log\epsilon}
        = \lim_{\epsilon\to 0}
            \frac{\log V(\epsilon)}{\log\epsilon} = \lambda
    \end{equation*}
    by Equation \eqref{eq:solve-LLC}
    
\end{solution}

\begin{solution}[ex:dln-lc]
\begin{enumerate}
    % Part (a)
    \item
        We first check which case of \cref{thm:aoyagi-watanabe} applies.
        Setting $M = N = H = m$, the three conditions of Case~1 each reduce to
        $m + r \leq 2m$, i.e.\ $r \leq m$, which always holds. So we are always
        in Case~1. The parity condition is on $M + H + N + r = 3m + r$, which
        has the same parity as $m + r$.

        \textbf{Case 1(a):} $m + r$ even. Setting $M = N = H = m$ in the
        formula:
        \begin{align*}
            \lambda
            &= \frac{-(m+r)^2 - m^2 - m^2 + 2(m+r)m + 2(m+r)m + 2m^2}{8} \\
            &= \frac{-(m+r)^2 + 4m(m+r)}{8} \\
            &= \frac{(m+r)\bigl(4m - (m+r)\bigr)}{8} \\
            &= \frac{(m+r)(3m-r)}{8},
        \end{align*}
        with $\mu = 1$.

        \textbf{Case 1(b):} $m + r$ odd. The same calculation gives
        \[
            \lambda = \frac{(m+r)(3m - r) + 1}{8},
        \]
        with $\mu = 2$.

    % Part (b)
    \item
         Differentiating with respect to $r$:
        \[
            \frac{d\lambda}{dr}=\frac{1}{8}\frac{d}{dr}\!\left[(m+r)(3m-r)\right]
            = \frac{1}{8}\left[(3m - r) - (m + r)\right] = \frac{1}{4}(m - r) \geq 0,
        \]
        so $\lambda$ is increasing in $r$ therefore \textbf{lower rank
        $\implies$ lower $\lambda$}. This is consistent with
        \cref{ex:dln-degeneracy:general-rank} where we saw that at lower-rank
        parameters, more directions in parameter space are degenerate.

    % Part (c)
    \item
        The parameter space is $\W = \R^{2m^2}$, so the regular-model baseline
        is $d/2 = m^2$. At full rank $r = m$:
        \[
            \lambda = \frac{(2m)(2m)}{8} = \frac{m^2}{2} = \frac{d}{4},
        \]
        which is already half the regular-model value. At rank $r = 0$:
        \[
            \lambda = \frac{m \cdot 3m}{8} = \frac{3m^2}{8} = \frac{3d}{16},
        \]
        even smaller. In fact $\lambda < d/2$ for every rank $0 \leq r \leq m$.

        This tells us that the two-layer DLN is singular at every point in its
        loss landscape, even at full-rank parameters.
\end{enumerate}
\end{solution}

% \clearpage
% \subsection[Solutions to exercises from Section~4]{Solutions to exercises from \cref{sec:fef}}
% \addtocontents{loe}{\solutionlistskip}

\begin{solution}[ex:fef-numerical]
    This exercise is computational. We outline the procedure and key
    observations.

    The model is $p(x \mid \mu) = (2\pi)^{-1/2}
    \exp\!\big({-}\tfrac{1}{2}(x - \mu^3)^2\big)$ with prior
    $\varphi(\mu) = (2\pi)^{-1/2}\exp(-\mu^2/2)$ and true distribution
    $q(x) = (2\pi)^{-1/2}\exp\!\big({-}\tfrac{1}{2}(x-\mu_0^3)^2\big)$.

    For each $n$, sample $x^{(1)}, \ldots, x^{(n)} \sim q$ and compute
    \begin{equation*}
        Z_n
        = \int_{-\infty}^{\infty}
            \varphi(\mu)
            \prod_{i=1}^n p(x^{(i)} \mid \mu)\, d\mu
    \end{equation*}
    the then set $F_n = -\log Z_n$.

    The empirical loss is
    \begin{equation*}
        L_n(\mu)
        = -\frac{1}{n}\sum_{i=1}^n \log p(x^{(i)} \mid \mu)
        = \frac{1}{2}\log(2\pi)
          + \frac{1}{2n}\sum_{i=1}^n (x^{(i)} - \mu^3)^2.
    \end{equation*}

    \begin{enumerate}
        \item
            For $\mu_0 = 0$, the LLC is $\lambda = 1/6$
            (\cref{ex:cubic-parameterisation}). Plot $F_n$ against $n$
            and compare with the asymptotic estimate
            $\hat F_n = n L_n(0) + \tfrac{1}{6}\log n$.
            For large $n$, the two curves should agree up to an
            $O(1)$ offset: $F_n - \hat F_n$ should stabilise
            around a constant.
        \item
            For $\mu_0 = 5$, the LLC is $\lambda = 1/2$ (regular case).
            The asymptotic estimate is
            $\hat F_n = n L_n(5) + \tfrac{1}{2}\log n$.
        \item
            For $\mu_0 = 0.1$, the asymptotic LLC is $\lambda = 1/2$
            (since $\mu_0 \neq 0$). However, the minimum of $L$ is at
            $\mu = \mu_0 = 0.1$, very close to the singularity at
            $\mu = 0$ where $\lambda = 1/6$. For moderate $n$, the
            posterior is spread broadly enough that it ``feels'' the
            nearby singularity, and $F_n$ behaves more like the
            $\lambda = 1/6$ regime from part~(a). Only for sufficiently
            large $n$ does the posterior concentrate tightly enough
            around $\mu_0 = 0.1$ for the regular $\lambda = 1/2$
            asymptotics to take hold.
            This demonstrates that the singularity at $\mu = 0$
            influences learning even when $\mu_0 \neq 0$, provided
            $\mu_0$ is close to the singular point. \qedhere
    \end{enumerate}
\end{solution}

\begin{solution}[ex:phase-transitions]
\begin{enumerate}
    \item
        By definition, the partition function is
        \begin{equation*}
            Z_n = \int_\W \varphi(w) \prod_{i=1}^n p(x^{(i)} \mid w)\, dw.
        \end{equation*}
        Since $\W = \U \sqcup \V$ is a disjoint union, the integral splits additively:
        \begin{equation*}
            Z_n = \int_\U \varphi(w) \prod_{i=1}^n p(x^{(i)} \mid w)\, dw
                + \int_\V \varphi(w) \prod_{i=1}^n p(x^{(i)} \mid w)\, dw
                = Z_n(\U) + Z_n(\V).
        \end{equation*}
        Recalling that $F_n = -\log Z_n$ and similarly for the local free energies,
        \begin{equation*}
            Z_n = e^{-F_n(\U)} + e^{-F_n(\V)},
        \end{equation*}
        and taking the negative logarithm gives
        \begin{equation*}
            F_n = -\log Z_n = -\log(e^{-F_n(\U)} + e^{-F_n(\V)}).
        \end{equation*}

    \item
        The curves $F_n$ and $\min(F_n(\U), F_n(\V))$ are nearly
        identical, both track whichever local free energy is smaller.
        However, $F_n$ is smooth whereas $\min(F_n(\U), F_n(\V))$ has a
        kink at $n \approx 570$, where $F_n(\U) = F_n(\V)$.
        For $n \lesssim 570$ the simpler region $\V$ (lower $\lambda$)
        dominates; for $n \gtrsim 570$ the lower loss region $\U$
        (lower loss coefficient) dominates.

    \item
        Differentiating the log-sum-exp from (a),
        \begin{equation*}
            \frac{d}{dn}F_n
            = \frac{\frac{d}{dn}F_n(\U)\,e^{-F_n(\U)} + \frac{d}{dn}F_n(\V)\,e^{-F_n(\V)}}
                   {e^{-F_n(\U)} + e^{-F_n(\V)}},
        \end{equation*}
        Plotting this reveals a steep step around
        $n \approx 570$, transitioning from $\frac{d}{dn}F_n(\V) \approx 0.5$ to
        $\frac{d}{dn}F_n(\U) \approx 0.3$. This rapid change in the derivative of
        $F_n$ is a smooth approximation of a discontinuity, justifying
        the term \emph{phase transition}.

    \item
        Write $F_n(\U) \approx L_\U \, n + \lambda_\U \log n$ and
        $F_n(\V) \approx L_\V \, n + \lambda_\V \log n$, and set
        $\Delta L = L_\U - L_\V$ and
        $\Delta\lambda = \lambda_\U - \lambda_\V$. The critical sample
        size $n^*$ where the phase transition occurs satisfies
        \[
            \Delta L \cdot n^* = -\Delta\lambda \cdot \log n^*,
            \qquad\text{i.e.}\qquad
            \frac{n^*}{\log n^*} = -\frac{\Delta\lambda}{\Delta L}
        \]
        Since $n^*/\log n^*$ is positive, a solution exists if
        and only if $-\Delta\lambda / \Delta L > 0$, i.e.\ $\Delta L$
        and $\Delta\lambda$ have \emph{opposite signs}. A phase transition occurs precisely when one region fits better
        while the other is simpler ($\Delta L$ and $\Delta\lambda$ have
        opposite signs). If both fit and complexity favour the same
        region, that region dominates for all $n$ and no phase transition
        occurs.

        $n^*/\log n^*$ is monotonically increasing for $n > e$ so for such $n$, increasing $|\Delta \lambda|$ increases $n^*$ and increasing $\Delta L$ decreases $n^*$.
        
        \qedhere
\end{enumerate}
\end{solution}

\end{document}
